\documentclass[12pt]{article}
\usepackage{amsmath}
\usepackage{amsfonts}
\usepackage{float}
\usepackage{fancyhdr}
\usepackage{graphicx}
\usepackage{setspace}
\usepackage[colorlinks=true,linkcolor=blue, citecolor=red]{hyperref}
\usepackage{url}
\usepackage[top=.75in, left=.75in, right=.75in, bottom=1in]{geometry}

% For algorithm
\usepackage{algorithm}
\usepackage{algpseudocode}

% ============ CODE ============
\usepackage{listings}
\usepackage{xcolor}
\definecolor{codegreen}{rgb}{0,0.6,0}
\definecolor{codegray}{rgb}{0.5,0.5,0.5}
\definecolor{codepurple}{rgb}{0.58,0,0.82}
\definecolor{backcolour}{rgb}{0.95,0.95,0.92}

% Styling for the code.
\lstdefinestyle{mystyle}{
    backgroundcolor=\color{backcolour},   
    commentstyle=\color{codegreen},
    keywordstyle=\color{magenta},
    numberstyle=\tiny\color{codegray},
    stringstyle=\color{codepurple},
    basicstyle=\ttfamily\footnotesize,
    breakatwhitespace=false,         
    breaklines=true,                 
    captionpos=b,                    
    keepspaces=true,                 
    numbers=left,                    
    numbersep=5pt,                  
    showspaces=false,                
    showstringspaces=false,
    showtabs=false,                  
    tabsize=2
}
\lstset{style=mystyle}

% Table packages and Unicode / Font support
\usepackage{array}
\usepackage{longtable}
\usepackage{enumitem}
\usepackage{booktabs}
\usepackage[utf8]{inputenc}
\DeclareUnicodeCharacter{2264}{\ensuremath{\le}}
\DeclareUnicodeCharacter{2194}{\ensuremath{\leftrightarrow}}
\usepackage[T1]{fontenc}
\usepackage{lmodern}
\graphicspath{{../../../../templates/}{../../prototype/}{figures/}{./}}
\setlength{\emergencystretch}{3em}
\setlist[enumerate]{label=\arabic*.,leftmargin=*}
\setlist[itemize]{leftmargin=*}

% ============ HEADER AND FOOTER ============
\setlength{\headheight}{29.43912pt}

\pagestyle{fancy}
\fancyfoot{}
\fancyfoot[R]{Page \thepage}

\setlength{\parindent}{0pt}
\renewcommand{\baselinestretch}{1.5}

\newcommand{\coursename}{Introduction of Software Engineering}
\newcommand{\reportname}{Studitopia}
\newcommand{\reporttitle}{\href{https://github.com/RoundFloor/Studitopia.git}{Studitopia}}

\newcommand{\studentname}{%
Nguyen Thanh An (24125024)\\
Nguyen Chi Bao (24125025)\\
Tran Nguyen Khanh Binh (24125026)\\
Tran Canh Anh Tuan (24125048)\\
Nguyen Le Quoc Huy (24125100)
}
\newcommand{\teachername}{Ho Tuan Thanh, Nguyen Thi Minh Tuyen}

% Header
\lhead{\reporttitle}
\rhead{
University of Science - VNUHCM\\
\coursename
}

% Footer
\newcommand{\leftfooter}{}
\lfoot{\leftfooter}

% ============ DOCUMENT ============
\begin{document}

\pagenumbering{roman}
\begin{titlepage}
\newcommand{\HRule}{\rule{\linewidth}{0.5mm}}
\centering

\textsc{\LARGE VNU-HCM}\\[0.5cm]
\textsc{\Large University of Science}\\[0.5cm]
\textsc{\large Faculty of Information Technology}\\[0.5cm]
\textsc{Computer Science}\\[0.5cm]

\HRule \\[0.4cm]
{ 
\huge{\bfseries{\reporttitle}}\\[0.5cm]
\large{\bfseries{Friends and Social Networking --- Use Case Diagram}}
}\\[0.4cm]
\HRule \\[0.5cm]

\textbf{\large Subject: \coursename}\\[0.5cm]

\begin{minipage}[t]{0.55\textwidth}
\begin{flushleft} \large
\emph{Students:}\\
\studentname
\end{flushleft}
\end{minipage}
~
\begin{minipage}[t]{0.4\textwidth}
\begin{flushright} \large
\emph{Supervising Teachers:} \\
\teachername
\end{flushright}
\end{minipage}\\[1cm]

{\large \today}\\[1cm]

\includegraphics[scale=.20]{img/hcmus-logo.png}\\[1cm] 

\vfill
\end{titlepage}	

\tableofcontents
\pagebreak

\pagenumbering{arabic}
\setcounter{page}{1}

\hypertarget{sec:friends-and-social-networking-use-case-diagram}{}
\addcontentsline{toc}{section}{Friends and Social Networking --- Use Case Diagram}
\section*{Friends and Social Networking --- Use Case Diagram}
\par\medskip

\textbf{Version:} \texttt{1.0}
\par\medskip

\textbf{Performed by:} [Nguyen Chi Bao] | \textbf{Edited by:} [Nguyen Chi Bao]
\par\medskip

Merged from the focused use case diagrams in this functional group.
\par\medskip

\begin{center}
\includegraphics[width=0.85\textwidth,height=0.75\textheight,keepaspectratio]{"figures/diagram_1.png"}
\end{center}
\par\medskip

\textbf{Diagram notation.} Person shape = human actor; cloud = supporting system; blue = use case; solid line = association; labelled solid arrow = generalization.
\par\medskip

\pagebreak
\hypertarget{sec:uc-26-find-a-study-partner}{}
\addcontentsline{toc}{section}{UC-26 --- Find a Study Partner}
\section*{UC-26 --- Find a Study Partner}
\par\medskip

\textbf{Version:} \texttt{1.0}
\par\medskip

\textbf{Performed by:} [Nguyen Chi Bao] | \textbf{Edited by:} [Nguyen Chi Bao]
\par\medskip

\begingroup
\small
\renewcommand{\baselinestretch}{1.0}\selectfont
\renewcommand{\arraystretch}{1.15}
\setlength{\tabcolsep}{6pt}
\begin{longtable}{|>{\raggedright\arraybackslash}p{0.09\textwidth}|>{\raggedright\arraybackslash}p{0.84\textwidth}|}
\hline
\textbf{Field} & \textbf{Value} \\ \hline
\endhead
Use Case ID & UC-26 \\ \hline
Name & Find a Study Partner \\ \hline
Primary Actor & University Student \\ \hline
Supporting Actors & None \\ \hline
Feature Group & Friends and Social Networking \\ \hline
Priority & Not explicitly stated (Not specified in proposals) \\ \hline
Related User Stories & US-30 \\ \hline
Related Requirements & \texttt{AC-30.1}, \texttt{AC-30.2}, \texttt{AC-30.3}, \texttt{FR-30.1}, \texttt{FR-30.2}, \texttt{FR-30.3}, \texttt{FR-30.4}, \texttt{AC-09.1}, \texttt{AC-09.2}, \texttt{AC-09.3}, \texttt{AC-09.4}, \texttt{AC-09.5}, \texttt{FR-09.1}, \texttt{FR-09.2}, \texttt{FR-09.3}, \texttt{FR-09.4}, \texttt{FR-09.5}, \texttt{AC-10.1}, \texttt{AC-10.2}, \texttt{AC-10.3}, \texttt{AC-10.4}, \texttt{FR-10.1}, \texttt{FR-10.2}, \texttt{FR-10.3}, \texttt{FR-10.4}, \texttt{NFR-PERF-1}, \texttt{NFR-SEC-1}, \texttt{NFR-SEC-3}, \texttt{NFR-SEC-6}, \texttt{NFR-PRIV-1}, \texttt{NFR-PRIV-2}, \texttt{NFR-PRIV-3}, \texttt{NFR-PRIV-4}, \texttt{NFR-PRIV-7}, \texttt{NFR-USE-2}, \texttt{NFR-USE-4}, \texttt{NFR-USE-5} \\ \hline
Source Documents & \texttt{proposals/userstory.md} (section 9.1); \texttt{proposals/functional\_requirement.md} (US-30); \texttt{proposals/friend.md}; \texttt{proposals/friend\_profile.md}; \texttt{proposals/group\_detail.md}; \texttt{proposals/settings.md}; \texttt{proposals/overview.md}; \texttt{proposals/feature\_flow.md}; \texttt{proposals/screen\_flow.md}; \texttt{proposals/API.md}; \texttt{proposals/response\_code.md}; \texttt{proposals/non\_functional\_requirement.md} \\ \hline
\end{longtable}
\endgroup
\par\medskip

\hypertarget{sec:part-a-use-case-model}{}
\addcontentsline{toc}{subsection}{Part A --- Use Case Model}
\subsection*{Part A --- Use Case Model}
\par\medskip

\hypertarget{sec:a1-focused-use-case-diagram}{}
\addcontentsline{toc}{subsubsection}{A1. Focused Use Case Diagram}
\subsubsection*{A1. Focused Use Case Diagram}
\par\medskip

\begin{center}
\includegraphics[width=0.85\textwidth,height=0.75\textheight,keepaspectratio]{"figures/UC-26-find-study-partner_1.png"}
\end{center}
\par\medskip

\textbf{Diagram notation.} Person shape = human actor; cloud = supporting system; blue = selected use case; solid line = association; labelled solid arrow = generalization; dashed arrow = include or extend.
\par\medskip

\hypertarget{sec:a2-actor-register}{}
\addcontentsline{toc}{subsubsection}{A2. Actor Register}
\subsubsection*{A2. Actor Register}
\par\medskip

\begingroup
\small
\renewcommand{\baselinestretch}{1.0}\selectfont
\renewcommand{\arraystretch}{1.15}
\setlength{\tabcolsep}{6pt}
\begin{longtable}{|>{\raggedright\arraybackslash}p{0.10\textwidth}|>{\raggedright\arraybackslash}p{0.09\textwidth}|>{\raggedright\arraybackslash}p{0.72\textwidth}|}
\hline
\textbf{Actor} & \textbf{Type} & \textbf{Description} \\ \hline
\endhead
University Student & Primary & The actor in US-30 who wants to find users with similar courses and learning goals and establish useful academic connections. \\ \hline
\end{longtable}
\endgroup
\par\medskip

\hypertarget{sec:a3-relationships}{}
\addcontentsline{toc}{subsubsection}{A3. Relationships}
\subsubsection*{A3. Relationships}
\par\medskip

\hypertarget{sec:a3-1-associations}{}
\paragraph*{A3.1 Associations}
\par\medskip

\begingroup
\small
\renewcommand{\baselinestretch}{1.0}\selectfont
\renewcommand{\arraystretch}{1.15}
\setlength{\tabcolsep}{6pt}
\begin{longtable}{|>{\raggedright\arraybackslash}p{0.12\textwidth}|>{\raggedright\arraybackslash}p{0.18\textwidth}|>{\raggedright\arraybackslash}p{0.09\textwidth}|>{\raggedright\arraybackslash}p{0.51\textwidth}|}
\hline
\textbf{Actor} & \textbf{Use Case} & \textbf{Type} & \textbf{Description} \\ \hline
\endhead
University Student & UC-26 Find a Study Partner & Primary & Searches and filters the member directory and reviews concise profile cards. \\ \hline
\end{longtable}
\endgroup
\par\medskip

\hypertarget{sec:a3-2-generalization-and-specialization}{}
\paragraph*{A3.2 Generalization and Specialization}
\par\medskip

\begingroup
\small
\renewcommand{\baselinestretch}{1.0}\selectfont
\renewcommand{\arraystretch}{1.15}
\setlength{\tabcolsep}{6pt}
\begin{longtable}{|>{\raggedright\arraybackslash}p{0.09\textwidth}|>{\raggedright\arraybackslash}p{0.09\textwidth}|>{\raggedright\arraybackslash}p{0.45\textwidth}|>{\raggedright\arraybackslash}p{0.29\textwidth}|}
\hline
\textbf{Child} & \textbf{Parent} & \textbf{Type} & \textbf{Inherited Behavior or Association} \\ \hline
\endhead
None & None & No actor or use-case generalization is named for US-30. & No inherited behavior is specified. \\ \hline
\end{longtable}
\endgroup
\par\medskip

\hypertarget{sec:a3-3-include}{}
\paragraph*{A3.3 Include}
\par\medskip

\begingroup
\small
\renewcommand{\baselinestretch}{1.0}\selectfont
\renewcommand{\arraystretch}{1.15}
\setlength{\tabcolsep}{6pt}
\begin{longtable}{|>{\raggedright\arraybackslash}p{0.14\textwidth}|>{\raggedright\arraybackslash}p{0.18\textwidth}|>{\raggedright\arraybackslash}p{0.60\textwidth}|}
\hline
\textbf{Base Use Case} & \textbf{Included Use Case} & \textbf{Meaning} \\ \hline
\endhead
None & None & No distinct included use case is named in the proposals. \\ \hline
\end{longtable}
\endgroup
\par\medskip

\hypertarget{sec:a3-4-extend}{}
\paragraph*{A3.4 Extend}
\par\medskip

\begingroup
\small
\renewcommand{\baselinestretch}{1.0}\selectfont
\renewcommand{\arraystretch}{1.15}
\setlength{\tabcolsep}{6pt}
\begin{longtable}{|>{\raggedright\arraybackslash}p{0.14\textwidth}|>{\raggedright\arraybackslash}p{0.10\textwidth}|>{\raggedright\arraybackslash}p{0.21\textwidth}|>{\raggedright\arraybackslash}p{0.45\textwidth}|}
\hline
\textbf{Extending Use Case} & \textbf{Base Use Case} & \textbf{Condition} & \textbf{Added Behavior} \\ \hline
\endhead
None & None & No condition is specified. & No distinct extending use case is named in the proposals. \\ \hline
\end{longtable}
\endgroup
\par\medskip

\hypertarget{sec:a4-model-notes}{}
\addcontentsline{toc}{subsubsection}{A4. Model Notes}
\subsubsection*{A4. Model Notes}
\par\medskip

\begin{enumerate}
\item The model contains only the selected use case and the story's University Student actor because no friend/group I-xx, E-xx, or G-xx use case is named in the proposals.
\item Recommended is inside the Friend Group area: \texttt{screen\_flow.md} gives \texttt{Home -$>$ Friend Group -$>$ Find friend}, and \texttt{feature\_flow.md} gives \texttt{Friend -$>$ Find friend -$>$ Friend item}.
\item \texttt{friend.md} recommends \texttt{/friends/find}, \texttt{StudyPartnerDiscoveryPanel}, \texttt{PartnerSearchInput}, and \texttt{FriendProfileCard}; these are UI recommendations, not API contracts.
\item Discovery is scoped by \texttt{FR-30.1}--\texttt{FR-30.4} and \texttt{AC-30.1}--\texttt{AC-30.3}; preference data comes from \texttt{FR-09.1}--\texttt{FR-09.5} and \texttt{AC-09.1}--\texttt{AC-09.5}.
\item Visibility of profile cards, mutual context, and shared forum context must follow \texttt{FR-10.1}--\texttt{FR-10.4}, \texttt{AC-10.1}--\texttt{AC-10.4}, and \texttt{NFR-PRIV-7}.
\end{enumerate}
\par\medskip

\hypertarget{sec:part-b-use-case-specification}{}
\addcontentsline{toc}{subsection}{Part B --- Use Case Specification}
\subsection*{Part B --- Use Case Specification}
\par\medskip

\hypertarget{sec:b1-brief-description}{}
\addcontentsline{toc}{subsubsection}{B1. Brief Description}
\subsubsection*{B1. Brief Description}
\par\medskip

A university student opens the Friend Group area and searches the member directory for peers with similar courses, subjects, learning goals, learning pace, and language. The system returns concise profile cards with mutual academic interests and, when available, mutual friends or shared forum contexts, so the student can decide whether to open a profile or start the separate friend-request flow.
\par\medskip

This use case covers discovery and evaluation only. Sending a request belongs to UC-27, and managing incoming requests belongs to UC-28. The forum remains the primary collaboration space; friend discovery supports it by turning shared academic context into durable connections.
\par\medskip

\hypertarget{sec:b2-preconditions}{}
\addcontentsline{toc}{subsubsection}{B2. Preconditions}
\subsubsection*{B2. Preconditions}
\par\medskip

\begin{enumerate}
\item The student is authenticated with an active session (\texttt{NFR-SEC-1}).
\item The student can reach the Friend Group from Home (\texttt{screen\_flow.md}).
\item Member-directory and learning-preference data are available to the system.
\item The student's own preferences may be used for matching (\texttt{FR-09.4}, \texttt{AC-09.4}).
\item Privacy settings and study-partner search opt-outs are enforced before results are displayed (\texttt{FR-10.2}, \texttt{AC-10.2}, \texttt{NFR-PRIV-7}).
\item Sensitive communication uses TLS 1.2 or higher (\texttt{NFR-SEC-3}).
\end{enumerate}
\par\medskip

\hypertarget{sec:b3-postconditions}{}
\addcontentsline{toc}{subsubsection}{B3. Postconditions}
\subsubsection*{B3. Postconditions}
\par\medskip

\begin{enumerate}
\item \textbf{Success:} The system displays matching profile cards with mutual academic interests (\texttt{FR-30.3}, \texttt{AC-30.2}).
\item \textbf{Success:} Mutual friends or shared forum contexts are displayed when available (\texttt{FR-30.4}, \texttt{AC-30.3}).
\item \textbf{Success:} Only members permitted by privacy settings appear in results (\texttt{FR-10.2}, \texttt{NFR-PRIV-7}).
\item \textbf{Failure:} The system shows an appropriate empty or restricted state and does not reveal hidden member data.
\end{enumerate}
\par\medskip

\hypertarget{sec:b4-main-success-scenario-basic-flow}{}
\addcontentsline{toc}{subsubsection}{B4. Main Success Scenario (Basic Flow)}
\subsubsection*{B4. Main Success Scenario (Basic Flow)}
\par\medskip

\begin{enumerate}
\item The student opens the Friend Group from Home.
\item The student selects \textbf{Find Friends} or the \texttt{FindFriendsTab}.
\item The system opens \texttt{StudyPartnerDiscoveryPanel} with \texttt{PartnerSearchInput}.
\item The student enters keywords using “Search by name, subject, course, or goal.”
\item The student selects filters for subject, goal, learning pace, and language (\texttt{FR-30.2}, \texttt{AC-30.1}).
\item The student may add the university, major, or course context using \texttt{PartnerAcademicContextFilter}.
\item The system applies privacy and search opt-out rules before computing results (\texttt{FR-10.2}, \texttt{NFR-PRIV-2}, \texttt{NFR-PRIV-7}).
\item The system displays \texttt{FriendProfileCard} items with university, major, shared courses, public hobby, mutual interests, and mutual connections (\texttt{FR-30.1}, \texttt{FR-30.3}, \texttt{FR-30.4}).
\item The system explains relevance with text such as “Shares CS101 and Database” or “3 mutual friends.”
\item The student opens a profile or proceeds to UC-27 Send Friend Request.
\end{enumerate}
\par\medskip

\hypertarget{sec:b5-alternative-flows}{}
\addcontentsline{toc}{subsubsection}{B5. Alternative Flows}
\subsubsection*{B5. Alternative Flows}
\par\medskip

\hypertarget{sec:a1-recommended-peers}{}
\paragraph*{A1. Recommended peers}
\par\medskip

At basic flow step 2, the student opens \textbf{Recommended}. \texttt{RecommendedStudyPartnerList} shows suggested peers with mutual interests and shared context (\texttt{feature\_flow.md}, \texttt{AC-30.2}, \texttt{AC-30.3}).
\par\medskip

\hypertarget{sec:a2-open-a-full-friend-profile}{}
\paragraph*{A2. Open a full friend profile}
\par\medskip

At basic flow step 10, the student selects \textbf{View profile}. \texttt{FriendProfilePage} shows the public academic context and mutual-context information the member permits (\texttt{friend\_profile.md}).
\par\medskip

\hypertarget{sec:a3-no-matching-study-partners}{}
\paragraph*{A3. No matching study partners}
\par\medskip

At basic flow step 8, no members match. The system shows “No study partners found” and suggests changing filters (\texttt{friend.md}, recommendation).
\par\medskip

\hypertarget{sec:a4-clear-filters}{}
\paragraph*{A4. Clear filters}
\par\medskip

At any point after step 5, the student selects \textbf{Clear filters}. Active chips are removed and discovery returns to the unfiltered state (\texttt{friend.md}, recommendation).
\par\medskip

\hypertarget{sec:b6-exception-flows}{}
\addcontentsline{toc}{subsubsection}{B6. Exception Flows}
\subsubsection*{B6. Exception Flows}
\par\medskip

No friend/group-specific error codes are specified; the following are source-backed UI states rather than backend error contracts.
\par\medskip

\hypertarget{sec:e1-student-opts-out-of-study-partner-search}{}
\paragraph*{E1. Student opts out of study-partner search}
\par\medskip

If another member has opted out, the system omits that member from results and does not expose a hidden profile (\texttt{FR-10.2}, \texttt{AC-10.2}).
\par\medskip

\hypertarget{sec:e2-hidden-fields}{}
\paragraph*{E2. Hidden fields}
\par\medskip

If a member hides academic or online-status information, the system omits those fields and may show “Some information is hidden by this member's privacy settings.” (\texttt{friend\_profile.md}, \texttt{NFR-PRIV-3}, \texttt{NFR-PRIV-4}).
\par\medskip

\hypertarget{sec:e3-discovery-data-fails-to-load}{}
\paragraph*{E3. Discovery data fails to load}
\par\medskip

The system shows “Unable to load friends or groups.” with Retry (\texttt{friend.md}, recommendation); no API status code is asserted.
\par\medskip

\hypertarget{sec:b7-business-rules}{}
\addcontentsline{toc}{subsubsection}{B7. Business Rules}
\subsubsection*{B7. Business Rules}
\par\medskip

\begin{enumerate}
\item The system must allow a user to find other university students with similar courses, subjects, or learning goals (\texttt{FR-30.1}).
\item Search and filtering must support subject, goal, learning pace, and language (\texttt{FR-30.2}, \texttt{AC-30.1}).
\item Results must display concise profile cards with mutual academic interests (\texttt{FR-30.3}, \texttt{AC-30.2}).
\item Mutual friends or shared forum contexts are displayed when available (\texttt{FR-30.4}, \texttt{AC-30.3}).
\item Preferences are used to recommend study partners (\texttt{FR-09.4}, \texttt{AC-09.4}).
\item Users can opt out of appearing in study-partner search results (\texttt{FR-10.2}, \texttt{AC-10.2}, \texttt{NFR-PRIV-2}).
\item Profile visibility, academic-info visibility, and online-status rules apply to discovery surfaces (\texttt{FR-10.1}, \texttt{FR-10.3}, \texttt{FR-10.4}, \texttt{NFR-PRIV-3}, \texttt{NFR-PRIV-4}, \texttt{NFR-PRIV-7}).
\item Protected resources require authentication, and authorization is enforced server-side (\texttt{NFR-SEC-1}, \texttt{NFR-SEC-6}).
\end{enumerate}
\par\medskip

\hypertarget{sec:b8-special-requirements}{}
\addcontentsline{toc}{subsubsection}{B8. Special Requirements}
\subsubsection*{B8. Special Requirements}
\par\medskip

\begingroup
\small
\renewcommand{\baselinestretch}{1.0}\selectfont
\renewcommand{\arraystretch}{1.15}
\setlength{\tabcolsep}{6pt}
\begin{longtable}{|>{\raggedright\arraybackslash}p{0.11\textwidth}|>{\raggedright\arraybackslash}p{0.82\textwidth}|}
\hline
\textbf{ID} & \textbf{Requirement} \\ \hline
\endhead
NFR-PERF-1 & Standard user interactions must respond within 2 seconds. \\ \hline
NFR-SEC-1 & Authentication is required before accessing protected user resources. \\ \hline
NFR-SEC-3 & Sensitive communication must use TLS 1.2 or higher. \\ \hline
NFR-SEC-6 & Server-side authorization checks must prevent unauthorized access. \\ \hline
NFR-PRIV-2 & Users must be able to allow or reject incoming friend requests. \\ \hline
NFR-PRIV-7 & Privacy settings must apply consistently to member lists and search results. \\ \hline
NFR-USE-2 & The interface must support responsive layouts from 320 px to 1920 px. \\ \hline
NFR-USE-4 & Primary actions must use clear labels and recognizable controls. \\ \hline
NFR-USE-5 & The interface must be usable at screen widths of at least 320 px. \\ \hline
\end{longtable}
\endgroup
\par\medskip

\hypertarget{sec:b9-extension-points}{}
\addcontentsline{toc}{subsubsection}{B9. Extension Points}
\subsubsection*{B9. Extension Points}
\par\medskip

\begingroup
\small
\renewcommand{\baselinestretch}{1.0}\selectfont
\renewcommand{\arraystretch}{1.15}
\setlength{\tabcolsep}{6pt}
\begin{longtable}{|>{\raggedright\arraybackslash}p{0.48\textwidth}|>{\raggedright\arraybackslash}p{0.45\textwidth}|}
\hline
\textbf{Extension Point} & \textbf{Location in Flow of Events} \\ \hline
\endhead
Recommended peers & Friend Group entry; alternative A1 \\ \hline
Opening a friend profile & Basic flow step 10; alternative A2 \\ \hline
Empty discovery state & Basic flow step 8; alternative A3 \\ \hline
Filter reset & Basic flow steps 5-8; alternative A4 \\ \hline
Separation to UC-27 Send Friend Request & Basic flow step 10 \\ \hline
\end{longtable}
\endgroup
\par\medskip

\hypertarget{sec:b10-interface-and-data-notes}{}
\addcontentsline{toc}{subsubsection}{B10. Interface and Data Notes}
\subsubsection*{B10. Interface and Data Notes}
\par\medskip

\begin{itemize}
\item \textbf{Screen/flow path (confirmed from proposals):} \texttt{screen\_flow.md} gives \texttt{Home -$>$ Friend Group -$>$ Find friend}; \texttt{feature\_flow.md} gives \texttt{Friend -$>$ Find friend -$>$ Friend item}; \texttt{userstory.md} section 9 gives \texttt{Friend -$>$ FindFriend -$>$ FriendItem}.
\item \textbf{Recommended routes and components:} \texttt{friend.md} recommends \texttt{/friends/find}, \texttt{FriendPage}, \texttt{FriendTabBar}, \texttt{FindFriendsTab}, \texttt{StudyPartnerDiscoveryPanel}, \texttt{PartnerSearchInput}, filter controls, \texttt{PartnerResultCountText}, and \texttt{PartnerNoResultsState}. These are UI recommendations, not specified requirements.
\item \textbf{Exact filter and field names:} subject, goal, learning pace, language, university, major, course, \texttt{displayName}, \texttt{avatarUrl}, \texttt{courses}, \texttt{subjects}, \texttt{learningGoals}, \texttt{learningPace}, \texttt{language}, \texttt{publicHobby}, \texttt{mutualInterests}, \texttt{mutualFriends}, and \texttt{sharedForumContexts} (\texttt{FR-30.1}--\texttt{FR-30.4}, \texttt{friend.md}, \texttt{friend\_profile.md}).
\item \textbf{Exact quoted UI text:} “Find study partners”, “Search by name, subject, course, or goal.”, “Clear filters”, “Shares CS101 and Database”, “3 mutual friends.”, “No study partners found”, “View profile”, and “Some information is hidden by this member's privacy settings.”
\item \textbf{Data rules:} \texttt{FR-30.3}/\texttt{AC-30.2} require mutual academic interests on cards; \texttt{FR-30.4}/\texttt{AC-30.3} require mutual friends or shared forum context when available. Hidden values must not leak through badges, counters, or supporting text (\texttt{NFR-PRIV-7}).
\item \textbf{API and status codes:} Not specified in proposals. No friend/search endpoint, HTTP status, or error code is defined for discovery.
\item \textbf{Notification/privacy rules:} This use case does not send a friend request or notification. Privacy notices and search opt-outs are confirmed requirements; component names and empty-state wording are recommendations.
\end{itemize}
\par\medskip

\hypertarget{sec:b11-source-alignment-notes}{}
\addcontentsline{toc}{subsubsection}{B11. Source Alignment Notes}
\subsubsection*{B11. Source Alignment Notes}
\par\medskip

\begin{itemize}
\item \texttt{use\_cases\_description.md} and the selected PA3 use-case set contain no dedicated friend/group UC; this specification is derived from US-30, the friend flows, and UI recommendations in \texttt{friend.md} and \texttt{friend\_profile.md}.
\item No friend/group API endpoints or error codes are defined in \texttt{API.md} or \texttt{response\_code.md}; no endpoint, HTTP status, or error code is invented here.
\item The source conflicts on whether the group owner or Superuser removes members: \texttt{FR-33.5}/\texttt{AC-33.4} name the group owner, while \texttt{feature\_flow.md:145} says Superuser kicks. This conflict is outside UC-26 but affects the same friend/group feature family.
\item \texttt{AC-30.1} names “subjects, goals, learning pace, and language”; \texttt{FR-30.2} names the singular equivalents. Both are preserved in the flow and business rules.
\item \texttt{friend.md} recommends example text such as “18 students found.”; it is example UI copy, not a required numeric result count.
\item Routes such as \texttt{/friends/find} and \texttt{/friends/users/:userId} are explicitly recommended in \texttt{friend.md} and are not API contracts.
\item UC-26 stops before the request action; duplicate prevention, the 200-character note limit, and request statuses are specified in UC-27 and UC-28.
\end{itemize}
\par\medskip

\par\medskip
\hypertarget{sec:ui-prototype-specimens}{}
\addcontentsline{toc}{subsection}{UI Prototype Specimens}
\subsection*{UI Prototype Specimens}
\par\medskip
\begin{figure}[H]
\centering
\includegraphics[width=0.75\textwidth,height=0.42\textheight,keepaspectratio]{"../../prototype/friend_01_friend_hub.png"}
\par\vspace{2pt}
{\footnotesize \textbf{UI Prototype:} Friend: 01 Friend hub}
\end{figure}
\par\smallskip
\begin{figure}[H]
\centering
\includegraphics[width=0.75\textwidth,height=0.42\textheight,keepaspectratio]{"../../prototype/friend_01_new_friend_daniel_park.png"}
\par\vspace{2pt}
{\footnotesize \textbf{UI Prototype:} Friend: 01 New friend - Daniel Park}
\end{figure}
\par\smallskip
\begin{figure}[H]
\centering
\includegraphics[width=0.75\textwidth,height=0.42\textheight,keepaspectratio]{"../../prototype/friend_01_new_friend_isabella_rossi.png"}
\par\vspace{2pt}
{\footnotesize \textbf{UI Prototype:} Friend: 01 New friend - Isabella Rossi}
\end{figure}
\par\smallskip
\begin{figure}[H]
\centering
\includegraphics[width=0.75\textwidth,height=0.42\textheight,keepaspectratio]{"../../prototype/friend_01_new_friend_ravi_shah.png"}
\par\vspace{2pt}
{\footnotesize \textbf{UI Prototype:} Friend: 01 New friend - Ravi Shah}
\end{figure}
\par\smallskip
\begin{figure}[H]
\centering
\includegraphics[width=0.75\textwidth,height=0.42\textheight,keepaspectratio]{"../../prototype/friend_02_find_friend_and_recommended.png"}
\par\vspace{2pt}
{\footnotesize \textbf{UI Prototype:} Friend: 02 Find friend and recommended}
\end{figure}
\par\smallskip
\begin{figure}[H]
\centering
\includegraphics[width=0.75\textwidth,height=0.42\textheight,keepaspectratio]{"../../prototype/friend_04_discovery_find_friend_2.png"}
\par\vspace{2pt}
{\footnotesize \textbf{UI Prototype:} Friend: 04 - Discovery --- Find friend (2)}
\end{figure}
\par\smallskip

\pagebreak
\hypertarget{sec:uc-27-send-friend-request}{}
\addcontentsline{toc}{section}{UC-27 --- Send Friend Request}
\section*{UC-27 --- Send Friend Request}
\par\medskip

\textbf{Version:} \texttt{1.0}
\par\medskip

\textbf{Performed by:} [Nguyen Chi Bao] | \textbf{Edited by:} [Nguyen Chi Bao]
\par\medskip

\begingroup
\small
\renewcommand{\baselinestretch}{1.0}\selectfont
\renewcommand{\arraystretch}{1.15}
\setlength{\tabcolsep}{6pt}
\begin{longtable}{|>{\raggedright\arraybackslash}p{0.09\textwidth}|>{\raggedright\arraybackslash}p{0.84\textwidth}|}
\hline
\textbf{Field} & \textbf{Value} \\ \hline
\endhead
Use Case ID & UC-27 \\ \hline
Name & Send Friend Request \\ \hline
Primary Actor & University Student \\ \hline
Supporting Actors & Friend Request Recipient; Notification Service \\ \hline
Feature Group & Friends and Social Networking \\ \hline
Priority & Not explicitly stated (Not specified in proposals) \\ \hline
Related User Stories & US-31 \\ \hline
Related Requirements & \texttt{AC-31.1}, \texttt{AC-31.2}, \texttt{AC-31.3}, \texttt{AC-31.4}, \texttt{FR-31.1}, \texttt{FR-31.2}, \texttt{FR-31.3}, \texttt{FR-31.4}, \texttt{FR-38.2}, \texttt{AC-38.3}, \texttt{FR-39.1}, \texttt{FR-39.3}, \texttt{FR-39.5}, \texttt{AC-39.1}, \texttt{AC-39.3}, \texttt{NFR-PRIV-2}, \texttt{NFR-PRIV-7}, \texttt{NFR-SEC-1}, \texttt{NFR-SEC-3}, \texttt{NFR-SEC-6} \\ \hline
Source Documents & \texttt{proposals/userstory.md} (section 9.2); \texttt{proposals/functional\_requirement.md} (US-31); \texttt{proposals/friend.md}; \texttt{proposals/friend\_profile.md}; \texttt{proposals/group\_detail.md}; \texttt{proposals/settings.md}; \texttt{proposals/overview.md}; \texttt{proposals/feature\_flow.md}; \texttt{proposals/screen\_flow.md}; \texttt{proposals/API.md}; \texttt{proposals/response\_code.md}; \texttt{proposals/non\_functional\_requirement.md} \\ \hline
\end{longtable}
\endgroup
\par\medskip

\hypertarget{sec:part-a-use-case-model}{}
\addcontentsline{toc}{subsection}{Part A --- Use Case Model}
\subsection*{Part A --- Use Case Model}
\par\medskip

\hypertarget{sec:a1-focused-use-case-diagram}{}
\addcontentsline{toc}{subsubsection}{A1. Focused Use Case Diagram}
\subsubsection*{A1. Focused Use Case Diagram}
\par\medskip

\begin{center}
\includegraphics[width=0.85\textwidth,height=0.75\textheight,keepaspectratio]{"figures/UC-27-send-friend-request_1.png"}
\end{center}
\par\medskip

\textbf{Diagram notation.} Person shape = human actor; cloud = supporting system; blue = selected use case; solid line = association; labelled solid arrow = generalization; dashed arrow = include or extend.
\par\medskip

\hypertarget{sec:a2-actor-register}{}
\addcontentsline{toc}{subsubsection}{A2. Actor Register}
\subsubsection*{A2. Actor Register}
\par\medskip

\begingroup
\small
\renewcommand{\baselinestretch}{1.0}\selectfont
\renewcommand{\arraystretch}{1.15}
\setlength{\tabcolsep}{6pt}
\begin{longtable}{|>{\raggedright\arraybackslash}p{0.15\textwidth}|>{\raggedright\arraybackslash}p{0.10\textwidth}|>{\raggedright\arraybackslash}p{0.67\textwidth}|}
\hline
\textbf{Actor} & \textbf{Type} & \textbf{Description} \\ \hline
\endhead
University Student & Primary & The US-31 actor who sends a friend request to maintain a connection for future learning activities. \\ \hline
Friend Request Recipient & Secondary & The other member who receives the request and may later accept, decline, or block under UC-28. \\ \hline
Notification Service & Supporting System & Delivers the social notification required by \texttt{FR-31.4} and \texttt{AC-31.4}; no notification endpoint is specified. \\ \hline
\end{longtable}
\endgroup
\par\medskip

\hypertarget{sec:a3-relationships}{}
\addcontentsline{toc}{subsubsection}{A3. Relationships}
\subsubsection*{A3. Relationships}
\par\medskip

\hypertarget{sec:a3-1-associations}{}
\paragraph*{A3.1 Associations}
\par\medskip

\begingroup
\small
\renewcommand{\baselinestretch}{1.0}\selectfont
\renewcommand{\arraystretch}{1.15}
\setlength{\tabcolsep}{6pt}
\begin{longtable}{|>{\raggedright\arraybackslash}p{0.15\textwidth}|>{\raggedright\arraybackslash}p{0.15\textwidth}|>{\raggedright\arraybackslash}p{0.10\textwidth}|>{\raggedright\arraybackslash}p{0.51\textwidth}|}
\hline
\textbf{Actor} & \textbf{Use Case} & \textbf{Type} & \textbf{Description} \\ \hline
\endhead
University Student & UC-27 Send Friend Request & Primary & Starts the request from a profile card or member list and optionally writes a note. \\ \hline
Friend Request Recipient & UC-27 Send Friend Request & Secondary & Receives the request and notification. \\ \hline
Notification Service & UC-27 Send Friend Request & Supporting & Routes the friend-request notification to the recipient. \\ \hline
\end{longtable}
\endgroup
\par\medskip

\hypertarget{sec:a3-2-generalization-and-specialization}{}
\paragraph*{A3.2 Generalization and Specialization}
\par\medskip

\begingroup
\small
\renewcommand{\baselinestretch}{1.0}\selectfont
\renewcommand{\arraystretch}{1.15}
\setlength{\tabcolsep}{6pt}
\begin{longtable}{|>{\raggedright\arraybackslash}p{0.12\textwidth}|>{\raggedright\arraybackslash}p{0.10\textwidth}|>{\raggedright\arraybackslash}p{0.10\textwidth}|>{\raggedright\arraybackslash}p{0.59\textwidth}|}
\hline
\textbf{Child} & \textbf{Parent} & \textbf{Type} & \textbf{Inherited Behavior or Association} \\ \hline
\endhead
Friend Request Recipient & University Student & Actor generalization & The recipient is another university student and inherits the general student association while adding the recipient role. \\ \hline
\end{longtable}
\endgroup
\par\medskip

\hypertarget{sec:a3-3-include}{}
\paragraph*{A3.3 Include}
\par\medskip

\begingroup
\small
\renewcommand{\baselinestretch}{1.0}\selectfont
\renewcommand{\arraystretch}{1.15}
\setlength{\tabcolsep}{6pt}
\begin{longtable}{|>{\raggedright\arraybackslash}p{0.14\textwidth}|>{\raggedright\arraybackslash}p{0.18\textwidth}|>{\raggedright\arraybackslash}p{0.60\textwidth}|}
\hline
\textbf{Base Use Case} & \textbf{Included Use Case} & \textbf{Meaning} \\ \hline
\endhead
None & None & No distinct included use case is named in the proposals. \\ \hline
\end{longtable}
\endgroup
\par\medskip

\hypertarget{sec:a3-4-extend}{}
\paragraph*{A3.4 Extend}
\par\medskip

\begingroup
\small
\renewcommand{\baselinestretch}{1.0}\selectfont
\renewcommand{\arraystretch}{1.15}
\setlength{\tabcolsep}{6pt}
\begin{longtable}{|>{\raggedright\arraybackslash}p{0.14\textwidth}|>{\raggedright\arraybackslash}p{0.10\textwidth}|>{\raggedright\arraybackslash}p{0.21\textwidth}|>{\raggedright\arraybackslash}p{0.45\textwidth}|}
\hline
\textbf{Extending Use Case} & \textbf{Base Use Case} & \textbf{Condition} & \textbf{Added Behavior} \\ \hline
\endhead
None & None & No condition is specified. & No distinct extending use case is named in the proposals. \\ \hline
\end{longtable}
\endgroup
\par\medskip

\hypertarget{sec:a4-model-notes}{}
\addcontentsline{toc}{subsubsection}{A4. Model Notes}
\subsubsection*{A4. Model Notes}
\par\medskip

\begin{enumerate}
\item The model contains only UC-27 and its actors because no friend/group I-xx, E-xx, or G-xx use case is named in the proposals.
\item \texttt{FR-31.1}/\texttt{AC-31.1} allow the request trigger from profile cards or member lists; \texttt{feature\_flow.md} names the stranger action \textbf{Make Friend}.
\item The recipient notification is a requirement, but \texttt{friend.md} and \texttt{friend\_profile.md} define no notification API and no friend error codes.
\item Request relationship states are \texttt{None}, \texttt{Pending}, \texttt{Accepted}, \texttt{Declined}, and \texttt{Blocked} (\texttt{friend\_profile.md}); UC-27 creates or prevents \texttt{Pending} and does not decide the outcome.
\item Duplicate prevention and the 200-character note limit are backend and frontend rules, not merely visual suggestions (\texttt{FR-31.2}, \texttt{FR-31.3}, \texttt{AC-31.2}, \texttt{AC-31.3}).
\end{enumerate}
\par\medskip

\hypertarget{sec:part-b-use-case-specification}{}
\addcontentsline{toc}{subsection}{Part B --- Use Case Specification}
\subsection*{Part B --- Use Case Specification}
\par\medskip

\hypertarget{sec:b1-brief-description}{}
\addcontentsline{toc}{subsubsection}{B1. Brief Description}
\subsubsection*{B1. Brief Description}
\par\medskip

A university student wants to connect with another member for future learning activities. The student selects \textbf{Add friend}, \textbf{Send friend request}, or the flow-derived \textbf{Make Friend} from a profile card, member list, or friend profile, optionally writes an introductory note of at most 200 characters, and submits the request. The system prevents duplicate requests and notifies the recipient.
\par\medskip

The recipient's decision is handled by UC-28 Manage Friend Requests. Successful sending changes the sender's view of the relationship to \texttt{Pending}; it does not create a confirmed friendship until the recipient accepts.
\par\medskip

\hypertarget{sec:b2-preconditions}{}
\addcontentsline{toc}{subsubsection}{B2. Preconditions}
\subsubsection*{B2. Preconditions}
\par\medskip

\begin{enumerate}
\item The sender is authenticated with an active session (\texttt{NFR-SEC-1}).
\item The sender can view the target profile card, member list entry, or friend profile.
\item The relationship state is \texttt{None}, not \texttt{Pending}, \texttt{Accepted}, \texttt{Declined}, or \texttt{Blocked}.
\item The recipient permits incoming friend requests (\texttt{NFR-PRIV-2}, \texttt{FR-39.1}, \texttt{FR-39.5}).
\item The recipient has not muted or blocked the sender from friend requests (\texttt{FR-39.3}, \texttt{AC-39.1}).
\item Sensitive communication uses TLS 1.2 or higher (\texttt{NFR-SEC-3}).
\end{enumerate}
\par\medskip

\hypertarget{sec:b3-postconditions}{}
\addcontentsline{toc}{subsubsection}{B3. Postconditions}
\subsubsection*{B3. Postconditions}
\par\medskip

\begin{enumerate}
\item \textbf{Success:} A friend request is stored with sender, recipient, note, sent time, and status \texttt{Pending}.
\item \textbf{Success:} The recipient is notified of the friend request (\texttt{FR-31.4}, \texttt{AC-31.4}, \texttt{FR-38.2}).
\item \textbf{Success:} The sender's relationship state shows \texttt{Pending}, and further send actions are blocked for that pair.
\item \textbf{Failure:} No request is stored, the status remains \texttt{None} or the existing status is unchanged, and the sender receives clear feedback.
\end{enumerate}
\par\medskip

\hypertarget{sec:b4-main-success-scenario-basic-flow}{}
\addcontentsline{toc}{subsubsection}{B4. Main Success Scenario (Basic Flow)}
\subsubsection*{B4. Main Success Scenario (Basic Flow)}
\par\medskip

\begin{enumerate}
\item The student opens a friend profile, profile card, or member list from the Friend Group area.
\item The system shows the relationship state \texttt{None} and an available \texttt{SendFriendRequestButton} or \texttt{MakeFriendButton}.
\item The student selects \textbf{Add friend}, \textbf{Send friend request}, or \textbf{Make Friend}.
\item The system opens \texttt{SendRequestDialog} with a recipient summary and \texttt{RequestNoteInput}.
\item The system displays the label \textbf{Message}, placeholder “Introduce yourself”, and counter \texttt{0/200}.
\item The student optionally types a note of at most 200 characters (\texttt{FR-31.3}, \texttt{AC-31.3}).
\item The system checks that no active request already exists between the pair (\texttt{FR-31.2}, \texttt{AC-31.2}).
\item The student selects \textbf{Send request}.
\item The system validates the note length and recipient permission, creates the request with status \texttt{Pending}, and records \texttt{sentAt}.
\item The Notification Service notifies the recipient (\texttt{FR-31.4}, \texttt{AC-31.4}, \texttt{FR-38.2}).
\item The sender sees \texttt{Pending} on \texttt{RequestStatusBadge}, and the request appears in \texttt{OutgoingRequestList}.
\end{enumerate}
\par\medskip

\hypertarget{sec:b5-alternative-flows}{}
\addcontentsline{toc}{subsubsection}{B5. Alternative Flows}
\subsubsection*{B5. Alternative Flows}
\par\medskip

\hypertarget{sec:a1-send-from-a-profile-card}{}
\paragraph*{A1. Send from a profile card}
\par\medskip

At basic flow step 1, the student uses \texttt{FriendProfileCard} in Find Friends or Recommended. The same dialog and validation apply (\texttt{AC-31.1}, \texttt{friend.md}).
\par\medskip

\hypertarget{sec:a2-send-from-a-friend-item-for-a-stranger}{}
\paragraph*{A2. Send from a friend item for a stranger}
\par\medskip

At basic flow step 3, the student uses the flow-derived \textbf{Make Friend} action from a stranger's friend item (\texttt{feature\_flow.md}).
\par\medskip

\hypertarget{sec:a3-cancel-composition}{}
\paragraph*{A3. Cancel composition}
\par\medskip

At basic flow step 8, the student selects \textbf{Cancel}. \texttt{SendRequestDialog} closes and no request is created (\texttt{friend.md}, recommendation).
\par\medskip

\hypertarget{sec:a4-existing-pending-request}{}
\paragraph*{A4. Existing pending request}
\par\medskip

At basic flow step 3 or 7, an active request already exists. The system shows “Friend request already sent.” and disables sending (\texttt{FR-31.2}, \texttt{AC-31.2}, \texttt{friend.md}, \texttt{friend\_profile.md}).
\par\medskip

\hypertarget{sec:a5-recipient-rejects-incoming-requests-by-setting}{}
\paragraph*{A5. Recipient rejects incoming requests by setting}
\par\medskip

At basic flow step 7, the recipient's communication setting rejects incoming friend requests. The send action is not offered or is denied (\texttt{NFR-PRIV-2}, \texttt{FR-39.5}).
\par\medskip

\hypertarget{sec:b6-exception-flows}{}
\addcontentsline{toc}{subsubsection}{B6. Exception Flows}
\subsubsection*{B6. Exception Flows}
\par\medskip

No friend/group-specific error codes are specified; the following are source-backed UI states rather than backend error contracts.
\par\medskip

\hypertarget{sec:e1-note-exceeds-200-characters}{}
\paragraph*{E1. Note exceeds 200 characters}
\par\medskip

The live counter exceeds \texttt{0/200} and the system blocks submission until the note is shortened (\texttt{AC-31.3}, \texttt{friend\_profile.md}).
\par\medskip

\hypertarget{sec:e2-duplicate-request-race}{}
\paragraph*{E2. Duplicate request race}
\par\medskip

The system does not create a second request and preserves the existing \texttt{Pending} state (\texttt{FR-31.2}, \texttt{AC-31.2}).
\par\medskip

\hypertarget{sec:e3-relationship-is-blocked}{}
\paragraph*{E3. Relationship is \texttt{Blocked}}
\par\medskip

The system shows “You cannot interact with this member.” and disables request controls (\texttt{friend\_profile.md}, \texttt{FR-39.3}, \texttt{FR-39.5}).
\par\medskip

\hypertarget{sec:e4-notification-delivery-failure}{}
\paragraph*{E4. Notification delivery failure}
\par\medskip

The request remains \texttt{Pending}; notification behavior is not rolled back or assigned a friend error code because none is defined.
\par\medskip

\hypertarget{sec:e5-profile-or-action-fails-to-load}{}
\paragraph*{E5. Profile or action fails to load}
\par\medskip

The system shows “Unable to load this profile.” with Retry (\texttt{friend\_profile.md}, recommendation); no endpoint or status code is asserted.
\par\medskip

\hypertarget{sec:b7-business-rules}{}
\addcontentsline{toc}{subsubsection}{B7. Business Rules}
\subsubsection*{B7. Business Rules}
\par\medskip

\begin{enumerate}
\item A user can send a friend request to another user (\texttt{FR-31.1}).
\item Requests can be triggered from profile cards or member lists (\texttt{AC-31.1}).
\item Duplicate requests between the same users are prevented (\texttt{FR-31.2}, \texttt{AC-31.2}).
\item The sender may attach a short introductory note (\texttt{FR-31.3}, \texttt{AC-31.3}).
\item The note is limited to 200 characters, represented by a \texttt{0/200} counter (\texttt{friend.md}, \texttt{friend\_profile.md}).
\item The recipient is notified of the friend request (\texttt{FR-31.4}, \texttt{AC-31.4}).
\item The request status values are \texttt{None}, \texttt{Pending}, \texttt{Accepted}, \texttt{Declined}, and \texttt{Blocked} (\texttt{friend\_profile.md}).
\item Incoming friend requests may be allowed or rejected by the recipient (\texttt{NFR-PRIV-2}, \texttt{FR-39.1}, \texttt{FR-39.5}).
\item Blocked interactions must be prevented server-side, not only by hiding controls (\texttt{NFR-SEC-6}, \texttt{FR-39.5}).
\end{enumerate}
\par\medskip

\hypertarget{sec:b8-special-requirements}{}
\addcontentsline{toc}{subsubsection}{B8. Special Requirements}
\subsubsection*{B8. Special Requirements}
\par\medskip

\begingroup
\small
\renewcommand{\baselinestretch}{1.0}\selectfont
\renewcommand{\arraystretch}{1.15}
\setlength{\tabcolsep}{6pt}
\begin{longtable}{|>{\raggedright\arraybackslash}p{0.12\textwidth}|>{\raggedright\arraybackslash}p{0.81\textwidth}|}
\hline
\textbf{ID} & \textbf{Requirement} \\ \hline
\endhead
NFR-SEC-1 & Authentication is required before accessing protected user resources. \\ \hline
NFR-SEC-3 & Sensitive communication must use TLS 1.2 or higher. \\ \hline
NFR-SEC-6 & Server-side authorization checks must prevent unauthorized access. \\ \hline
NFR-PRIV-2 & Users must be able to allow or reject incoming friend requests. \\ \hline
NFR-PRIV-7 & Privacy settings must apply consistently to member lists and search results. \\ \hline
Design constraint & The request note must accept at most 200 characters, and duplicate requests must be prevented for the same pair. \\ \hline
\end{longtable}
\endgroup
\par\medskip

\hypertarget{sec:b9-extension-points}{}
\addcontentsline{toc}{subsubsection}{B9. Extension Points}
\subsubsection*{B9. Extension Points}
\par\medskip

\begingroup
\small
\renewcommand{\baselinestretch}{1.0}\selectfont
\renewcommand{\arraystretch}{1.15}
\setlength{\tabcolsep}{6pt}
\begin{longtable}{|>{\raggedright\arraybackslash}p{0.33\textwidth}|>{\raggedright\arraybackslash}p{0.60\textwidth}|}
\hline
\textbf{Extension Point} & \textbf{Location in Flow of Events} \\ \hline
\endhead
Profile-card request trigger & Basic flow step 1; alternative A1 \\ \hline
Flow-derived \textbf{Make Friend} trigger & Basic flow step 3; alternative A2 \\ \hline
Duplicate request state & Basic flow steps 3 and 7; alternative A4; exception E2 \\ \hline
Recipient permission check & Basic flow steps 7-9; alternative A5 \\ \hline
Notification delivery & Basic flow step 10; exception E4 \\ \hline
Handoff to UC-28 & After basic flow step 11, when the recipient manages the request \\ \hline
\end{longtable}
\endgroup
\par\medskip

\hypertarget{sec:b10-interface-and-data-notes}{}
\addcontentsline{toc}{subsubsection}{B10. Interface and Data Notes}
\subsubsection*{B10. Interface and Data Notes}
\par\medskip

\begin{itemize}
\item \textbf{Screen/flow path (confirmed from proposals):} \texttt{screen\_flow.md} gives \texttt{Home -$>$ Friend Group -$>$ Friend Request}; \texttt{feature\_flow.md} gives \texttt{Friend -$>$ FriendRequest -$>$ FriendItem} and \texttt{FriendItem -$>$ Stranger -$>$ Make Friend}; \texttt{friend.md} recommends \texttt{FindFriendsTab}, \texttt{FriendProfilePage}, and \texttt{SendRequestDialog}.
\item \textbf{Recommended routes and components:} \texttt{/friends/users/:userId}, \texttt{/friends/users/:userId/request}, \texttt{SendFriendRequestButton}, \texttt{RequestNoteInput}, \texttt{SubmitFriendRequestButton}, \texttt{CancelFriendRequestButton}, \texttt{DuplicateRequestState}, \texttt{RequestStatusBadge}, and \texttt{OutgoingRequestList} are UI recommendations from \texttt{friend.md} and \texttt{friend\_profile.md}, not confirmed API contracts.
\item \textbf{Exact request data fields:} \texttt{requestId}, \texttt{sender}, \texttt{recipient}, \texttt{note}, \texttt{status}, \texttt{sentAt}, and \texttt{notificationState} (\texttt{friend.md}); \texttt{friend\_profile.md} uses \texttt{requestNote} for the note field. Confirmed statuses are \texttt{None}, \texttt{Pending}, \texttt{Accepted}, \texttt{Declined}, and \texttt{Blocked}.
\item \textbf{Exact quoted UI text:} \textbf{Message}, “Introduce yourself”, \texttt{0/200}, \textbf{Add friend}, \textbf{Send friend request}, \textbf{Make Friend}, \textbf{Send request}, \textbf{Cancel}, “Friend request already sent.”, “Friend request accepted/declined.”, “Your friend request was accepted/declined.”, and “You cannot interact with this member.”
\item \textbf{Privacy notices:} “Some information is hidden by this member's privacy settings.” and “Direct messages may be limited to confirmed friends.” are source-backed notices on the viewed profile (\texttt{friend\_profile.md}).
\item \textbf{API endpoints, statuses, and error codes:} Not specified in proposals. No friend-request route, HTTP status, or error code is asserted; UI error banners are recommendations.
\item \textbf{Notification/privacy rules:} Recipient notification is a confirmed requirement (\texttt{FR-31.4}, \texttt{AC-31.4}, \texttt{FR-38.2}). Recipient opt-out, mute, and block are confirmed controls (\texttt{NFR-PRIV-2}, \texttt{FR-39.3}, \texttt{FR-39.5}); badge and toast wording is recommended UI copy.
\end{itemize}
\par\medskip

\hypertarget{sec:b11-source-alignment-notes}{}
\addcontentsline{toc}{subsubsection}{B11. Source Alignment Notes}
\subsubsection*{B11. Source Alignment Notes}
\par\medskip

\begin{itemize}
\item There is no dedicated PA3 friend/group UC; UC-27 is derived from US-31 and the friend/profile UI specifications rather than from an existing UC file.
\item No friend/group API endpoints or error codes are defined in \texttt{API.md} or \texttt{response\_code.md}; duplicate prevention and notification are requirements without a specified backend contract.
\item The source conflicts on whether the group owner or Superuser removes members: \texttt{FR-33.5}/\texttt{AC-33.4} name the group owner, while \texttt{feature\_flow.md:145} says Superuser kicks. This conflict is outside UC-27 but affects the same friend/group feature family.
\item \texttt{friend.md} lists \texttt{Pending}, \texttt{Accepted}, \texttt{Declined}, and \texttt{Blocked}; \texttt{friend\_profile.md} adds \texttt{None}. All five exact statuses are preserved here.
\item \texttt{friend.md} says “Your friend request was accepted/declined.” while \texttt{friend\_profile.md} says “Friend request accepted/declined.” Both are retained as distinct quoted UI text.
\item \texttt{FR-31.3}/\texttt{AC-31.3} state the short note and 200-character limit; component names, placeholder, and counter format are recommendations that make the limit visible.
\item \texttt{FR-38.2}/\texttt{AC-38.3} cover social notifications, but no notification endpoint or delivery guarantee is specified for friend requests.
\item \texttt{settings.md} defines \texttt{FriendRequestPermissionSelect} with \textbf{Allow / Reject incoming friend requests}; this is a UI grouping of \texttt{FR-39.1}/\texttt{FR-39.5} and \texttt{NFR-PRIV-2}, not an additional API rule.
\end{itemize}
\par\medskip

\par\medskip
\hypertarget{sec:ui-prototype-specimens}{}
\addcontentsline{toc}{subsection}{UI Prototype Specimens}
\subsection*{UI Prototype Specimens}
\par\medskip
\begin{figure}[H]
\centering
\includegraphics[width=0.75\textwidth,height=0.42\textheight,keepaspectratio]{"../../prototype/friend_01_new_friend_daniel_park.png"}
\par\vspace{2pt}
{\footnotesize \textbf{UI Prototype:} Friend: 01 New friend - Daniel Park}
\end{figure}
\par\smallskip
\begin{figure}[H]
\centering
\includegraphics[width=0.75\textwidth,height=0.42\textheight,keepaspectratio]{"../../prototype/friend_01_new_friend_isabella_rossi.png"}
\par\vspace{2pt}
{\footnotesize \textbf{UI Prototype:} Friend: 01 New friend - Isabella Rossi}
\end{figure}
\par\smallskip
\begin{figure}[H]
\centering
\includegraphics[width=0.75\textwidth,height=0.42\textheight,keepaspectratio]{"../../prototype/friend_01_new_friend_ravi_shah.png"}
\par\vspace{2pt}
{\footnotesize \textbf{UI Prototype:} Friend: 01 New friend - Ravi Shah}
\end{figure}
\par\smallskip
\begin{figure}[H]
\centering
\includegraphics[width=0.75\textwidth,height=0.42\textheight,keepaspectratio]{"../../prototype/friend_03_friend_request.png"}
\par\vspace{2pt}
{\footnotesize \textbf{UI Prototype:} Friend: 03 Friend request}
\end{figure}
\par\smallskip

\pagebreak
\hypertarget{sec:uc-28-manage-friend-requests}{}
\addcontentsline{toc}{section}{UC-28 --- Manage Friend Requests}
\section*{UC-28 --- Manage Friend Requests}
\par\medskip

\textbf{Version:} \texttt{1.0}
\par\medskip

\textbf{Performed by:} [Nguyen Chi Bao] | \textbf{Edited by:} [Nguyen Chi Bao]
\par\medskip

\begingroup
\small
\renewcommand{\baselinestretch}{1.0}\selectfont
\renewcommand{\arraystretch}{1.15}
\setlength{\tabcolsep}{6pt}
\begin{longtable}{|>{\raggedright\arraybackslash}p{0.09\textwidth}|>{\raggedright\arraybackslash}p{0.84\textwidth}|}
\hline
\textbf{Field} & \textbf{Value} \\ \hline
\endhead
Use Case ID & UC-28 \\ \hline
Name & Manage Friend Requests \\ \hline
Primary Actor & User \\ \hline
Supporting Actors & Friend Request Sender; Notification Service \\ \hline
Feature Group & Friends and Social Networking \\ \hline
Priority & Not explicitly stated (Not specified in proposals) \\ \hline
Related User Stories & US-32 \\ \hline
Related Requirements & \texttt{AC-32.1}, \texttt{AC-32.2}, \texttt{AC-32.3}, \texttt{FR-32.1}, \texttt{FR-32.2}, \texttt{FR-32.3}, \texttt{FR-32.4}, \texttt{FR-32.5}, \texttt{AC-31.4}, \texttt{FR-31.4}, \texttt{FR-38.2}, \texttt{AC-38.3}, \texttt{FR-39.1}, \texttt{FR-39.3}, \texttt{FR-39.5}, \texttt{AC-39.1}, \texttt{AC-39.3}, \texttt{NFR-PRIV-2}, \texttt{NFR-PRIV-7}, \texttt{NFR-SEC-1}, \texttt{NFR-SEC-3}, \texttt{NFR-SEC-6} \\ \hline
Source Documents & \texttt{proposals/userstory.md} (section 9.3); \texttt{proposals/functional\_requirement.md} (US-32); \texttt{proposals/friend.md}; \texttt{proposals/friend\_profile.md}; \texttt{proposals/group\_detail.md}; \texttt{proposals/settings.md}; \texttt{proposals/overview.md}; \texttt{proposals/feature\_flow.md}; \texttt{proposals/screen\_flow.md}; \texttt{proposals/API.md}; \texttt{proposals/response\_code.md}; \texttt{proposals/non\_functional\_requirement.md} \\ \hline
\end{longtable}
\endgroup
\par\medskip

\hypertarget{sec:part-a-use-case-model}{}
\addcontentsline{toc}{subsection}{Part A --- Use Case Model}
\subsection*{Part A --- Use Case Model}
\par\medskip

\hypertarget{sec:a1-focused-use-case-diagram}{}
\addcontentsline{toc}{subsubsection}{A1. Focused Use Case Diagram}
\subsubsection*{A1. Focused Use Case Diagram}
\par\medskip

\begin{center}
\includegraphics[width=0.85\textwidth,height=0.75\textheight,keepaspectratio]{"figures/UC-28-manage-friend-requests_1.png"}
\end{center}
\par\medskip

\textbf{Diagram notation.} Person shape = human actor; cloud = supporting system; blue = selected use case; solid line = association; labelled solid arrow = generalization; dashed arrow = include or extend.
\par\medskip

\hypertarget{sec:a2-actor-register}{}
\addcontentsline{toc}{subsubsection}{A2. Actor Register}
\subsubsection*{A2. Actor Register}
\par\medskip

\begingroup
\small
\renewcommand{\baselinestretch}{1.0}\selectfont
\renewcommand{\arraystretch}{1.15}
\setlength{\tabcolsep}{6pt}
\begin{longtable}{|>{\raggedright\arraybackslash}p{0.12\textwidth}|>{\raggedright\arraybackslash}p{0.10\textwidth}|>{\raggedright\arraybackslash}p{0.70\textwidth}|}
\hline
\textbf{Actor} & \textbf{Type} & \textbf{Description} \\ \hline
\endhead
User & Primary & The US-32 actor who accepts or rejects incoming friend requests to control social connections. \\ \hline
Friend Request Sender & Secondary & The member who sent the request, receives the outcome notification, and may be added to the friend list. \\ \hline
Notification Service & Supporting System & Routes the acceptance or decline notification required by \texttt{FR-32.5} and \texttt{FR-31.4}; no notification endpoint is specified. \\ \hline
\end{longtable}
\endgroup
\par\medskip

\hypertarget{sec:a3-relationships}{}
\addcontentsline{toc}{subsubsection}{A3. Relationships}
\subsubsection*{A3. Relationships}
\par\medskip

\hypertarget{sec:a3-1-associations}{}
\paragraph*{A3.1 Associations}
\par\medskip

\begingroup
\small
\renewcommand{\baselinestretch}{1.0}\selectfont
\renewcommand{\arraystretch}{1.15}
\setlength{\tabcolsep}{6pt}
\begin{longtable}{|>{\raggedright\arraybackslash}p{0.14\textwidth}|>{\raggedright\arraybackslash}p{0.18\textwidth}|>{\raggedright\arraybackslash}p{0.10\textwidth}|>{\raggedright\arraybackslash}p{0.49\textwidth}|}
\hline
\textbf{Actor} & \textbf{Use Case} & \textbf{Type} & \textbf{Description} \\ \hline
\endhead
User & UC-28 Manage Friend Requests & Primary & Reviews incoming requests and accepts, declines, or blocks future requests. \\ \hline
Friend Request Sender & UC-28 Manage Friend Requests & Secondary & Is added to the friend list on acceptance or notified on decline. \\ \hline
Notification Service & UC-28 Manage Friend Requests & Supporting & Delivers the accepted or declined outcome to the sender. \\ \hline
\end{longtable}
\endgroup
\par\medskip

\hypertarget{sec:a3-2-generalization-and-specialization}{}
\paragraph*{A3.2 Generalization and Specialization}
\par\medskip

\begingroup
\small
\renewcommand{\baselinestretch}{1.0}\selectfont
\renewcommand{\arraystretch}{1.15}
\setlength{\tabcolsep}{6pt}
\begin{longtable}{|>{\raggedright\arraybackslash}p{0.12\textwidth}|>{\raggedright\arraybackslash}p{0.09\textwidth}|>{\raggedright\arraybackslash}p{0.12\textwidth}|>{\raggedright\arraybackslash}p{0.57\textwidth}|}
\hline
\textbf{Child} & \textbf{Parent} & \textbf{Type} & \textbf{Inherited Behavior or Association} \\ \hline
\endhead
Friend Request Sender & User & Actor generalization & The sender is also a user and inherits the general user association while adding the sender role. \\ \hline
\end{longtable}
\endgroup
\par\medskip

\hypertarget{sec:a3-3-include}{}
\paragraph*{A3.3 Include}
\par\medskip

\begingroup
\small
\renewcommand{\baselinestretch}{1.0}\selectfont
\renewcommand{\arraystretch}{1.15}
\setlength{\tabcolsep}{6pt}
\begin{longtable}{|>{\raggedright\arraybackslash}p{0.14\textwidth}|>{\raggedright\arraybackslash}p{0.18\textwidth}|>{\raggedright\arraybackslash}p{0.60\textwidth}|}
\hline
\textbf{Base Use Case} & \textbf{Included Use Case} & \textbf{Meaning} \\ \hline
\endhead
None & None & No distinct included use case is named in the proposals. \\ \hline
\end{longtable}
\endgroup
\par\medskip

\hypertarget{sec:a3-4-extend}{}
\paragraph*{A3.4 Extend}
\par\medskip

\begingroup
\small
\renewcommand{\baselinestretch}{1.0}\selectfont
\renewcommand{\arraystretch}{1.15}
\setlength{\tabcolsep}{6pt}
\begin{longtable}{|>{\raggedright\arraybackslash}p{0.14\textwidth}|>{\raggedright\arraybackslash}p{0.10\textwidth}|>{\raggedright\arraybackslash}p{0.21\textwidth}|>{\raggedright\arraybackslash}p{0.45\textwidth}|}
\hline
\textbf{Extending Use Case} & \textbf{Base Use Case} & \textbf{Condition} & \textbf{Added Behavior} \\ \hline
\endhead
None & None & No condition is specified. & No distinct extending use case is named in the proposals. \\ \hline
\end{longtable}
\endgroup
\par\medskip

\hypertarget{sec:a4-model-notes}{}
\addcontentsline{toc}{subsubsection}{A4. Model Notes}
\subsubsection*{A4. Model Notes}
\par\medskip

\begin{enumerate}
\item The model contains only UC-28 and its actors because no friend/group I-xx, E-xx, or G-xx use case is named in the proposals.
\item \texttt{screen\_flow.md} gives \texttt{Friend Group -$>$ Friend Request}; \texttt{feature\_flow.md} gives \texttt{Friend -$>$ FriendRequest -$>$ FriendItem}.
\item \texttt{friend.md} recommends \texttt{FriendRequestCenter}, \texttt{IncomingRequestList}, \texttt{OutgoingRequestList}, \texttt{AcceptRequestButton}, \texttt{RejectRequestButton}, and \texttt{BlockFutureRequestsToggle}; these are UI recommendations.
\item Request states are \texttt{None}, \texttt{Pending}, \texttt{Accepted}, \texttt{Declined}, and \texttt{Blocked}; UC-28 changes \texttt{Pending} to one of the three outcome states.
\item Accepting adds the sender to the friend list (\texttt{FR-32.2}, \texttt{AC-32.2}); declining removes the pending request (\texttt{FR-32.3}, \texttt{AC-32.2}).
\item Blocking future requests is a decline-path option (\texttt{FR-32.4}, \texttt{AC-32.3}) and is also reflected in communication settings (\texttt{settings.md}).
\end{enumerate}
\par\medskip

\hypertarget{sec:part-b-use-case-specification}{}
\addcontentsline{toc}{subsection}{Part B --- Use Case Specification}
\subsection*{Part B --- Use Case Specification}
\par\medskip

\hypertarget{sec:b1-brief-description}{}
\addcontentsline{toc}{subsubsection}{B1. Brief Description}
\subsubsection*{B1. Brief Description}
\par\medskip

A user opens a centralized request center to review incoming friend requests with profile details and notes. The user can accept a request, which adds the sender to the friend list, or decline it, which removes the pending request. When declining, the user can block future requests from that sender, and the sender is notified of the outcome.
\par\medskip

UC-28 begins after UC-27 has created a \texttt{Pending} request. It does not cover sending requests, direct messaging, or group membership.
\par\medskip

\hypertarget{sec:b2-preconditions}{}
\addcontentsline{toc}{subsubsection}{B2. Preconditions}
\subsubsection*{B2. Preconditions}
\par\medskip

\begin{enumerate}
\item The user is authenticated with an active session (\texttt{NFR-SEC-1}).
\item At least one incoming request with status \texttt{Pending} is available, or the request center shows an empty state.
\item The user can view the sender's profile details allowed by privacy settings (\texttt{FR-10.4}, \texttt{NFR-PRIV-7}).
\item The user permits incoming friend requests or has an existing pending request to decide (\texttt{NFR-PRIV-2}).
\item Sensitive communication uses TLS 1.2 or higher (\texttt{NFR-SEC-3}).
\end{enumerate}
\par\medskip

\hypertarget{sec:b3-postconditions}{}
\addcontentsline{toc}{subsubsection}{B3. Postconditions}
\subsubsection*{B3. Postconditions}
\par\medskip

\begin{enumerate}
\item \textbf{Success:} On acceptance, the sender is added to the friend list and the request status becomes \texttt{Accepted}.
\item \textbf{Success:} On decline, the pending request is removed or shown as \texttt{Declined} and the status becomes \texttt{Declined}.
\item \textbf{Success:} If the user selected the block option, future requests from the sender are blocked and the status is \texttt{Blocked}.
\item \textbf{Success:} The sender is notified of the accepted or declined outcome (\texttt{FR-32.5}, \texttt{FR-31.4}).
\item \textbf{Failure:} The request remains \texttt{Pending}, no friend connection is created, and the user receives clear feedback.
\end{enumerate}
\par\medskip

\hypertarget{sec:b4-main-success-scenario-basic-flow}{}
\addcontentsline{toc}{subsubsection}{B4. Main Success Scenario (Basic Flow)}
\subsubsection*{B4. Main Success Scenario (Basic Flow)}
\par\medskip

\begin{enumerate}
\item The user opens the Friend Group from Home.
\item The user selects the \textbf{Friend Requests} tab or \texttt{FriendRequestsTab}.
\item The system opens \texttt{FriendRequestCenter} with \texttt{IncomingRequestList} and \texttt{OutgoingRequestList}.
\item The system displays each incoming request with sender avatar/name, academic context, note, sent time, and status \texttt{Pending}.
\item The user reviews the request note of at most 200 characters.
\item The user selects \textbf{Accept}.
\item The system verifies that the request is still \texttt{Pending} and that the user is authorized to decide it (\texttt{NFR-SEC-6}).
\item The system adds the sender to the friend list and changes the status to \texttt{Accepted} (\texttt{FR-32.2}, \texttt{AC-32.2}).
\item The Notification Service notifies the sender of acceptance (\texttt{FR-32.5}, \texttt{FR-31.4}).
\item The user sees confirmation feedback, and the sender's profile now offers friend actions instead of a pending request.
\end{enumerate}
\par\medskip

\hypertarget{sec:b5-alternative-flows}{}
\addcontentsline{toc}{subsubsection}{B5. Alternative Flows}
\subsubsection*{B5. Alternative Flows}
\par\medskip

\hypertarget{sec:a1-decline-a-request}{}
\paragraph*{A1. Decline a request}
\par\medskip

At basic flow step 6, the user selects \textbf{Decline}. The system removes the pending request, sets the status to \texttt{Declined}, and notifies the sender (\texttt{FR-32.3}, \texttt{AC-32.2}, \texttt{FR-32.5}).
\par\medskip

\hypertarget{sec:a2-decline-and-block-future-requests}{}
\paragraph*{A2. Decline and block future requests}
\par\medskip

At basic flow step 6, the user enables “Block future requests from this user.” before selecting \textbf{Decline}. The request becomes \texttt{Blocked}, and future requests from that sender are prevented (\texttt{FR-32.4}, \texttt{AC-32.3}).
\par\medskip

\hypertarget{sec:a3-manage-from-the-sender-s-profile}{}
\paragraph*{A3. Manage from the sender's profile}
\par\medskip

At basic flow step 4, the user opens \texttt{FriendProfilePage} from the request card. \texttt{AcceptRequestButton} and \texttt{RejectRequestButton} perform the same decision with the same status and notification rules (\texttt{friend\_profile.md}).
\par\medskip

\hypertarget{sec:a4-no-incoming-requests}{}
\paragraph*{A4. No incoming requests}
\par\medskip

At basic flow step 4, the list is empty. The system shows an appropriate empty state and keeps the request center available.
\par\medskip

\hypertarget{sec:a5-mark-request-notification-as-read}{}
\paragraph*{A5. Mark request notification as read}
\par\medskip

After a decision, the user selects \textbf{Mark as read} on the social notification (\texttt{FR-38.6}, \texttt{settings.md}).
\par\medskip

\hypertarget{sec:b6-exception-flows}{}
\addcontentsline{toc}{subsubsection}{B6. Exception Flows}
\subsubsection*{B6. Exception Flows}
\par\medskip

No friend/group-specific error codes are specified; the following are source-backed UI states rather than backend error contracts.
\par\medskip

\hypertarget{sec:e1-stale-or-already-decided-request}{}
\paragraph*{E1. Stale or already decided request}
\par\medskip

The request is no longer \texttt{Pending}. The system does not add a duplicate friend connection and refreshes the request list.
\par\medskip

\hypertarget{sec:e2-unauthorized-decision-attempt}{}
\paragraph*{E2. Unauthorized decision attempt}
\par\medskip

The backend denies the action through server-side authorization checks and does not change the request status (\texttt{NFR-SEC-6}, \texttt{FR-39.5}).
\par\medskip

\hypertarget{sec:e3-notification-delivery-failure}{}
\paragraph*{E3. Notification delivery failure}
\par\medskip

The accepted, declined, or blocked state remains in effect; no friend error code is assigned because none is defined.
\par\medskip

\hypertarget{sec:e4-request-center-fails-to-load}{}
\paragraph*{E4. Request center fails to load}
\par\medskip

The system shows “Unable to load friends or groups.” with Retry (\texttt{friend.md}, recommendation); no endpoint or status code is asserted.
\par\medskip

\hypertarget{sec:e5-sender-is-blocked-or-restricted}{}
\paragraph*{E5. Sender is blocked or restricted}
\par\medskip

The system preserves the \texttt{Blocked} state and shows “You cannot interact with this member.” where interaction controls would appear (\texttt{friend\_profile.md}).
\par\medskip

\hypertarget{sec:b7-business-rules}{}
\addcontentsline{toc}{subsubsection}{B7. Business Rules}
\subsubsection*{B7. Business Rules}
\par\medskip

\begin{enumerate}
\item A centralized view must display incoming friend requests (\texttt{FR-32.1}, \texttt{AC-32.1}).
\item Accepting adds the sender to the friend list (\texttt{FR-32.2}, \texttt{AC-32.2}).
\item Declining removes the pending request (\texttt{FR-32.3}, \texttt{AC-32.2}).
\item When declining, the user can block future requests from the sender (\texttt{FR-32.4}, \texttt{AC-32.3}).
\item The sender is notified when a request is accepted or rejected (\texttt{FR-32.5}).
\item The exact request statuses are \texttt{None}, \texttt{Pending}, \texttt{Accepted}, \texttt{Declined}, and \texttt{Blocked} (\texttt{friend\_profile.md}).
\item Request notes are limited to 200 characters (\texttt{AC-31.3}, \texttt{friend.md}, \texttt{friend\_profile.md}).
\item Users can allow or reject incoming friend requests (\texttt{NFR-PRIV-2}) and mute or block specific members (\texttt{FR-39.3}, \texttt{AC-39.1}).
\item Communication and privacy settings must prevent interactions the user does not permit (\texttt{FR-39.5}, \texttt{NFR-SEC-6}).
\end{enumerate}
\par\medskip

\hypertarget{sec:b8-special-requirements}{}
\addcontentsline{toc}{subsubsection}{B8. Special Requirements}
\subsubsection*{B8. Special Requirements}
\par\medskip

\begingroup
\small
\renewcommand{\baselinestretch}{1.0}\selectfont
\renewcommand{\arraystretch}{1.15}
\setlength{\tabcolsep}{6pt}
\begin{longtable}{|>{\raggedright\arraybackslash}p{0.10\textwidth}|>{\raggedright\arraybackslash}p{0.83\textwidth}|}
\hline
\textbf{ID} & \textbf{Requirement} \\ \hline
\endhead
NFR-SEC-1 & Authentication is required before accessing protected user resources. \\ \hline
NFR-SEC-3 & Sensitive communication must use TLS 1.2 or higher. \\ \hline
NFR-SEC-6 & Server-side authorization checks must prevent unauthorized access. \\ \hline
NFR-PRIV-2 & Users must be able to allow or reject incoming friend requests. \\ \hline
NFR-PRIV-7 & Privacy settings must apply consistently to member lists and search results. \\ \hline
Design constraint & A decision must be applied once to a \texttt{Pending} request; duplicate acceptance, duplicate friend connections, and error-code-derived messages are not defined. \\ \hline
\end{longtable}
\endgroup
\par\medskip

\hypertarget{sec:b9-extension-points}{}
\addcontentsline{toc}{subsubsection}{B9. Extension Points}
\subsubsection*{B9. Extension Points}
\par\medskip

\begingroup
\small
\renewcommand{\baselinestretch}{1.0}\selectfont
\renewcommand{\arraystretch}{1.15}
\setlength{\tabcolsep}{6pt}
\begin{longtable}{|>{\raggedright\arraybackslash}p{0.31\textwidth}|>{\raggedright\arraybackslash}p{0.62\textwidth}|}
\hline
\textbf{Extension Point} & \textbf{Location in Flow of Events} \\ \hline
\endhead
Request center entry & Basic flow steps 1-3 \\ \hline
Accept action & Basic flow step 6; alternative A3 \\ \hline
Decline action & Alternative A1 \\ \hline
Block future requests & Alternative A2; exception E5 \\ \hline
Outcome notification & Basic flow step 9; alternative A1; exception E3 \\ \hline
Notification read state & Alternative A5 \\ \hline
\end{longtable}
\endgroup
\par\medskip

\hypertarget{sec:b10-interface-and-data-notes}{}
\addcontentsline{toc}{subsubsection}{B10. Interface and Data Notes}
\subsubsection*{B10. Interface and Data Notes}
\par\medskip

\begin{itemize}
\item \textbf{Screen/flow path (confirmed from proposals):} \texttt{screen\_flow.md} gives \texttt{Home -$>$ Friend Group -$>$ Friend Request}; \texttt{feature\_flow.md} gives \texttt{Friend -$>$ FriendRequest -$>$ FriendItem}; \texttt{friend.md} recommends \texttt{FriendRequestsTab}, \texttt{FriendRequestCenter}, \texttt{IncomingRequestList}, and \texttt{OutgoingRequestList}.
\item \textbf{Recommended components:} \texttt{FriendRequestCard}, \texttt{RequestNoteText}, \texttt{AcceptRequestButton}, \texttt{RejectRequestButton}, \texttt{BlockFutureRequestsToggle}, \texttt{RequestStatusBadge}, \texttt{RequestNotificationBadge}, and \texttt{FriendToastHost} are UI recommendations from \texttt{friend.md} and \texttt{friend\_profile.md}.
\item \textbf{Exact request fields and statuses:} \texttt{requestId}, \texttt{sender}, \texttt{recipient}, \texttt{note}, \texttt{status}, \texttt{sentAt}, \texttt{notificationState}; exact statuses are \texttt{None}, \texttt{Pending}, \texttt{Accepted}, \texttt{Declined}, and \texttt{Blocked}.
\item \textbf{Exact quoted UI text:} \textbf{Friend Requests}, \textbf{Accept}, \textbf{Decline}, “Block future requests from this user.”, “Friend request accepted/declined.”, “Your friend request was accepted/declined.”, “Friend request already sent.”, \textbf{Mark as read}, and “Unable to load friends or groups.” with Retry.
\item \textbf{Request note display:} \texttt{RequestNoteText} shows the introductory note with a maximum of 200 characters (\texttt{FR-31.3}, \texttt{AC-31.3}). Hidden profile fields must not leak into request cards (\texttt{NFR-PRIV-7}).
\item \textbf{Settings linkage:} \texttt{FriendRequestPermissionSelect} offers \textbf{Allow / Reject incoming friend requests}; \texttt{MutedUsersList} and \texttt{BlockedUsersList} reflect communication controls (\texttt{settings.md}, \texttt{FR-39.1}, \texttt{FR-39.3}, \texttt{FR-39.5}).
\item \textbf{API endpoints, statuses, and error codes:} Not specified in proposals. No friend-request decision route, HTTP status, or error code is asserted.
\item \textbf{Notification/privacy rules:} Sender notification is a confirmed requirement (\texttt{FR-32.5}, \texttt{FR-31.4}, \texttt{FR-38.2}); badge, toast, and empty-state wording is recommended UI copy.
\end{itemize}
\par\medskip

\hypertarget{sec:b11-source-alignment-notes}{}
\addcontentsline{toc}{subsubsection}{B11. Source Alignment Notes}
\subsubsection*{B11. Source Alignment Notes}
\par\medskip

\begin{itemize}
\item There is no dedicated PA3 friend/group UC; UC-28 is derived from US-32 and the request-center recommendations in \texttt{friend.md} and \texttt{friend\_profile.md}.
\item No friend/group API endpoints or error codes are defined in \texttt{API.md} or \texttt{response\_code.md}; accept, decline, and block have requirements but no specified backend contract.
\item The source conflicts on whether the group owner or Superuser removes members: \texttt{FR-33.5}/\texttt{AC-33.4} name the group owner, while \texttt{feature\_flow.md:145} says Superuser kicks. This conflict is outside UC-28 but affects the same friend/group feature family.
\item \texttt{AC-32.2} groups accept and decline together, while \texttt{FR-32.2} and \texttt{FR-32.3} state them separately; the flow separates them while preserving both exact requirements.
\item \texttt{friend.md} lists \texttt{Pending}, \texttt{Accepted}, \texttt{Declined}, and \texttt{Blocked}; \texttt{friend\_profile.md} adds \texttt{None}. All five exact statuses are preserved here.
\item \texttt{FR-32.5} requires sender notification for acceptance or rejection, but no notification endpoint or delivery guarantee is specified in \texttt{API.md}.
\item \texttt{settings.md} adds \textbf{Allow / Reject incoming friend requests} as a communication-setting UI grouping of \texttt{FR-39.1}/\texttt{FR-39.5} and \texttt{NFR-PRIV-2}; it does not define a separate decision API.
\item UC-28 does not cancel outgoing requests; \texttt{CancelOutgoingRequestButton} in \texttt{friend.md} is explicitly a recommendation.
\end{itemize}
\par\medskip

\par\medskip
\hypertarget{sec:ui-prototype-specimens}{}
\addcontentsline{toc}{subsection}{UI Prototype Specimens}
\subsection*{UI Prototype Specimens}
\par\medskip
\begin{figure}[H]
\centering
\includegraphics[width=0.75\textwidth,height=0.42\textheight,keepaspectratio]{"../../prototype/friend_03_friend_request.png"}
\par\vspace{2pt}
{\footnotesize \textbf{UI Prototype:} Friend: 03 Friend request}
\end{figure}
\par\smallskip

\pagebreak
\hypertarget{sec:uc-29-create-and-manage-friend-groups}{}
\addcontentsline{toc}{section}{UC-29 --- Create and Manage Friend Groups}
\section*{UC-29 --- Create and Manage Friend Groups}
\par\medskip

\textbf{Version:} \texttt{1.0}
\par\medskip

\textbf{Performed by:} [Nguyen Chi Bao] | \textbf{Edited by:} [Nguyen Chi Bao]
\par\medskip

\begingroup
\small
\renewcommand{\baselinestretch}{1.0}\selectfont
\renewcommand{\arraystretch}{1.15}
\setlength{\tabcolsep}{6pt}
\begin{longtable}{|>{\raggedright\arraybackslash}p{0.09\textwidth}|>{\raggedright\arraybackslash}p{0.84\textwidth}|}
\hline
\textbf{Field} & \textbf{Value} \\ \hline
\endhead
Use Case ID & UC-29 \\ \hline
Name & Create and Manage Friend Groups \\ \hline
Primary Actor & University Student \\ \hline
Supporting Actors & Group Owner; Group Member; Notification Service \\ \hline
Feature Group & Friends and Social Networking \\ \hline
Priority & Not explicitly stated (Not specified in proposals) \\ \hline
Related User Stories & US-33 \\ \hline
Related Requirements & \texttt{AC-33.1}, \texttt{AC-33.2}, \texttt{AC-33.3}, \texttt{AC-33.4}, \texttt{AC-33.5}, \texttt{FR-33.1}, \texttt{FR-33.2}, \texttt{FR-33.3}, \texttt{FR-33.4}, \texttt{FR-33.5}, \texttt{FR-33.6}, \texttt{AC-27.1}, \texttt{AC-27.2}, \texttt{AC-27.3}, \texttt{FR-27.1}, \texttt{FR-27.2}, \texttt{FR-27.3}, \texttt{FR-27.4}, \texttt{FR-27.5}, \texttt{FR-15.3}, \texttt{FR-15.5}, \texttt{FR-15.6}, \texttt{FR-10.4}, \texttt{AC-10.4}, \texttt{FR-38.2}, \texttt{AC-38.3}, \texttt{FR-39.1}, \texttt{FR-39.3}, \texttt{FR-39.5}, \texttt{AC-39.1}, \texttt{AC-39.3}, \texttt{FR-40.1}, \texttt{FR-40.2}, \texttt{FR-40.3}, \texttt{FR-40.4}, \texttt{AC-40.1}, \texttt{AC-40.2}, \texttt{AC-40.3}, \texttt{AC-40.4}, \texttt{NFR-SEC-1}, \texttt{NFR-SEC-3}, \texttt{NFR-SEC-4}, \texttt{NFR-SEC-5}, \texttt{NFR-SEC-6}, \texttt{NFR-SEC-8}, \texttt{NFR-PRIV-5}, \texttt{NFR-PRIV-6}, \texttt{NFR-PRIV-7}, \texttt{NFR-USE-2}, \texttt{NFR-USE-4}, \texttt{NFR-USE-5} \\ \hline
Source Documents & \texttt{proposals/userstory.md} (section 9.4); \texttt{proposals/functional\_requirement.md} (US-33); \texttt{proposals/friend.md}; \texttt{proposals/friend\_profile.md}; \texttt{proposals/group\_detail.md}; \texttt{proposals/settings.md}; \texttt{proposals/overview.md}; \texttt{proposals/feature\_flow.md}; \texttt{proposals/screen\_flow.md}; \texttt{proposals/API.md}; \texttt{proposals/response\_code.md}; \texttt{proposals/non\_functional\_requirement.md} \\ \hline
\end{longtable}
\endgroup
\par\medskip

\hypertarget{sec:part-a-use-case-model}{}
\addcontentsline{toc}{subsection}{Part A --- Use Case Model}
\subsection*{Part A --- Use Case Model}
\par\medskip

\hypertarget{sec:a1-focused-use-case-diagram}{}
\addcontentsline{toc}{subsubsection}{A1. Focused Use Case Diagram}
\subsubsection*{A1. Focused Use Case Diagram}
\par\medskip

\begin{center}
\includegraphics[width=0.85\textwidth,height=0.75\textheight,keepaspectratio]{"figures/UC-29-manage-friend-groups_1.png"}
\end{center}
\par\medskip

\textbf{Diagram notation.} Person shape = human actor; cloud = supporting system; blue = selected use case; solid line = association; labelled solid arrow = generalization; dashed arrow = include or extend.
\par\medskip

\hypertarget{sec:a2-actor-register}{}
\addcontentsline{toc}{subsubsection}{A2. Actor Register}
\subsubsection*{A2. Actor Register}
\par\medskip

\begingroup
\small
\renewcommand{\baselinestretch}{1.0}\selectfont
\renewcommand{\arraystretch}{1.15}
\setlength{\tabcolsep}{6pt}
\begin{longtable}{|>{\raggedright\arraybackslash}p{0.11\textwidth}|>{\raggedright\arraybackslash}p{0.10\textwidth}|>{\raggedright\arraybackslash}p{0.71\textwidth}|}
\hline
\textbf{Actor} & \textbf{Type} & \textbf{Description} \\ \hline
\endhead
University Student & Primary & The US-33 actor who creates and manages a private or public friend group for classmates and project teams. \\ \hline
Group Owner & Secondary & A student with the \texttt{Owner} role who creates the group, may delete it, and is named by \texttt{FR-33.5}/\texttt{AC-33.4} for member removal. \\ \hline
Group Member & Secondary & A student with the \texttt{Member} role who can view members, invite others, message, share documents, and leave. \\ \hline
Notification Service & Supporting System & Routes group-invitation and social-activity notifications from \texttt{FR-38.2}; no notification endpoint is specified. \\ \hline
\end{longtable}
\endgroup
\par\medskip

\hypertarget{sec:a3-relationships}{}
\addcontentsline{toc}{subsubsection}{A3. Relationships}
\subsubsection*{A3. Relationships}
\par\medskip

\hypertarget{sec:a3-1-associations}{}
\paragraph*{A3.1 Associations}
\par\medskip

\begingroup
\small
\renewcommand{\baselinestretch}{1.0}\selectfont
\renewcommand{\arraystretch}{1.15}
\setlength{\tabcolsep}{6pt}
\begin{longtable}{|>{\raggedright\arraybackslash}p{0.12\textwidth}|>{\raggedright\arraybackslash}p{0.22\textwidth}|>{\raggedright\arraybackslash}p{0.10\textwidth}|>{\raggedright\arraybackslash}p{0.47\textwidth}|}
\hline
\textbf{Actor} & \textbf{Use Case} & \textbf{Type} & \textbf{Description} \\ \hline
\endhead
University Student & UC-29 Create and Manage Friend Groups & Primary & Creates a group or manages participation in a group. \\ \hline
Group Owner & UC-29 Create and Manage Friend Groups & Secondary & Configures the group and may remove members under the source-backed owner rule. \\ \hline
Group Member & UC-29 Create and Manage Friend Groups & Secondary & Uses group chat, documents, member list, invitations, and leave action. \\ \hline
Notification Service & UC-29 Create and Manage Friend Groups & Supporting & Routes invitation and social-activity notifications. \\ \hline
\end{longtable}
\endgroup
\par\medskip

\hypertarget{sec:a3-2-generalization-and-specialization}{}
\paragraph*{A3.2 Generalization and Specialization}
\par\medskip

\begingroup
\small
\renewcommand{\baselinestretch}{1.0}\selectfont
\renewcommand{\arraystretch}{1.15}
\setlength{\tabcolsep}{6pt}
\begin{longtable}{|>{\raggedright\arraybackslash}p{0.10\textwidth}|>{\raggedright\arraybackslash}p{0.15\textwidth}|>{\raggedright\arraybackslash}p{0.17\textwidth}|>{\raggedright\arraybackslash}p{0.49\textwidth}|}
\hline
\textbf{Child} & \textbf{Parent} & \textbf{Type} & \textbf{Inherited Behavior or Association} \\ \hline
\endhead
Group Owner & University Student & Actor generalization & Inherits the student association and adds the \texttt{Owner} role. \\ \hline
Group Member & University Student & Actor generalization & Inherits the student association and adds the \texttt{Member} role. \\ \hline
\end{longtable}
\endgroup
\par\medskip

\hypertarget{sec:a3-3-include}{}
\paragraph*{A3.3 Include}
\par\medskip

\begingroup
\small
\renewcommand{\baselinestretch}{1.0}\selectfont
\renewcommand{\arraystretch}{1.15}
\setlength{\tabcolsep}{6pt}
\begin{longtable}{|>{\raggedright\arraybackslash}p{0.14\textwidth}|>{\raggedright\arraybackslash}p{0.18\textwidth}|>{\raggedright\arraybackslash}p{0.60\textwidth}|}
\hline
\textbf{Base Use Case} & \textbf{Included Use Case} & \textbf{Meaning} \\ \hline
\endhead
None & None & No distinct included use case is named in the proposals. \\ \hline
\end{longtable}
\endgroup
\par\medskip

\hypertarget{sec:a3-4-extend}{}
\paragraph*{A3.4 Extend}
\par\medskip

\begingroup
\small
\renewcommand{\baselinestretch}{1.0}\selectfont
\renewcommand{\arraystretch}{1.15}
\setlength{\tabcolsep}{6pt}
\begin{longtable}{|>{\raggedright\arraybackslash}p{0.14\textwidth}|>{\raggedright\arraybackslash}p{0.10\textwidth}|>{\raggedright\arraybackslash}p{0.21\textwidth}|>{\raggedright\arraybackslash}p{0.45\textwidth}|}
\hline
\textbf{Extending Use Case} & \textbf{Base Use Case} & \textbf{Condition} & \textbf{Added Behavior} \\ \hline
\endhead
None & None & No condition is specified. & No distinct extending use case is named in the proposals. \\ \hline
\end{longtable}
\endgroup
\par\medskip

\hypertarget{sec:a4-model-notes}{}
\addcontentsline{toc}{subsubsection}{A4. Model Notes}
\subsubsection*{A4. Model Notes}
\par\medskip

\begin{enumerate}
\item The model contains only UC-29 and its actors because no friend/group I-xx, E-xx, or G-xx use case is named in the proposals.
\item \texttt{screen\_flow.md} gives \texttt{Friend Group -$>$ Create Group}; \texttt{feature\_flow.md} gives \texttt{Friend -$>$ Create Group -$>$ Privacy / Beginning member / Name}.
\item \texttt{group\_detail.md} recommends \texttt{/friends/groups/:groupId}, \texttt{GroupDetailPage}, \texttt{GroupMessagePanel}, \texttt{SharedDocumentPanel}, and \texttt{GroupMemberPanel}; these are UI recommendations, not API contracts.
\item Exact roles are \texttt{Owner}, \texttt{Superuser}, and \texttt{Member}. The source conflict between owner removal and Superuser kick is recorded in B11 and is not silently resolved.
\item Privacy values are only \texttt{Private} and \texttt{Public}; no additional privacy level, capacity limit, invite expiry, or approval rule is specified.
\item Meeting-room rules are not copied into groups; the forum remains the primary collaboration space.
\end{enumerate}
\par\medskip

\hypertarget{sec:part-b-use-case-specification}{}
\addcontentsline{toc}{subsection}{Part B --- Use Case Specification}
\subsection*{Part B --- Use Case Specification}
\par\medskip

\hypertarget{sec:b1-brief-description}{}
\addcontentsline{toc}{subsubsection}{B1. Brief Description}
\subsubsection*{B1. Brief Description}
\par\medskip

A university student creates a private or public friend group by naming it, setting privacy, and choosing beginning members. Members then collaborate in \texttt{GroupDetailPage} by sending messages, sharing documents, viewing the member list, inviting new members, removing members where authorized, and leaving the group.
\par\medskip

Group deletion is overview-derived from \texttt{overview.md} and owner-scoped as a recommendation. Group invitations have no specified approval step, expiration period, or automatic acceptance.
\par\medskip

\hypertarget{sec:b2-preconditions}{}
\addcontentsline{toc}{subsubsection}{B2. Preconditions}
\subsubsection*{B2. Preconditions}
\par\medskip

\begin{enumerate}
\item The student is authenticated with an active session (\texttt{NFR-SEC-1}).
\item The student can reach the Friend Group from Home (\texttt{screen\_flow.md}).
\item For creation, the student supplies a group name and selects \texttt{Private} or \texttt{Public} (\texttt{FR-33.1}, \texttt{FR-33.2}, \texttt{AC-33.1}).
\item For management, the student is an authorized member of the group (\texttt{NFR-SEC-4}, \texttt{NFR-SEC-6}).
\item Invitees are available in the friend/user picker; no invite approval or expiry rule is assumed.
\item Sensitive communication uses TLS 1.2 or higher (\texttt{NFR-SEC-3}).
\end{enumerate}
\par\medskip

\hypertarget{sec:b3-postconditions}{}
\addcontentsline{toc}{subsubsection}{B3. Postconditions}
\subsubsection*{B3. Postconditions}
\par\medskip

\begin{enumerate}
\item \textbf{Success:} A group exists with a name, \texttt{Private} or \texttt{Public} privacy, an owner, and initial members (\texttt{FR-33.1}, \texttt{FR-33.2}, \texttt{AC-33.1}).
\item \textbf{Success:} Members can send messages and share documents within the group (\texttt{FR-33.3}, \texttt{AC-33.2}).
\item \textbf{Success:} Members can view the member list and invite new members (\texttt{FR-33.4}, \texttt{AC-33.3}).
\item \textbf{Success:} An authorized role can remove a member, and any member can leave (\texttt{FR-33.5}, \texttt{FR-33.6}, \texttt{AC-33.4}, \texttt{AC-33.5}).
\item \textbf{Failure:} The group is not created or changed, private content is not exposed, and the user receives clear feedback without an invented error code.
\end{enumerate}
\par\medskip

\hypertarget{sec:b4-main-success-scenario-basic-flow}{}
\addcontentsline{toc}{subsubsection}{B4. Main Success Scenario (Basic Flow)}
\subsubsection*{B4. Main Success Scenario (Basic Flow)}
\par\medskip

\begin{enumerate}
\item The student opens the Friend Group and selects \textbf{Create Group} (\texttt{screen\_flow.md}).
\item The system opens \texttt{CreateGroupDialog} with \texttt{GroupNameInput}, optional \texttt{GroupDescriptionInput}, \texttt{GroupPrivacySelect}, and \texttt{BeginningMemberPicker}.
\item The student enters a group name and selects \texttt{Private} or \texttt{Public} (\texttt{FR-33.1}, \texttt{FR-33.2}, \texttt{AC-33.1}).
\item The student chooses beginning members from friends/users.
\item The student selects \textbf{Create group}.
\item The system validates the setup, creates the group with the student as \texttt{Owner}, and adds the chosen beginning members.
\item The new group appears in \texttt{GroupList} and \texttt{MyGroupsPanel}.
\item The student opens \texttt{GroupDetailPage}, which shows \texttt{GroupHeader}, \texttt{GroupPrivacyBadge}, \texttt{GroupMemberCountText}, chat, documents, and members.
\item The student sends a message with \texttt{GroupMessageInput} and \textbf{Send} (\texttt{FR-33.3}, \texttt{AC-33.2}).
\item The student uploads a document with \textbf{Upload document}; the system validates file type and size before acceptance (\texttt{FR-27.3}, \texttt{FR-15.3}, \texttt{NFR-SEC-5}).
\item The student opens \texttt{GroupMemberPanel} and views the member list with \texttt{Owner}, \texttt{Superuser}, or \texttt{Member} role badges (\texttt{FR-33.4}, \texttt{AC-33.3}).
\item The system enforces access so private-group content is visible to authorized members only (\texttt{NFR-SEC-4}, \texttt{NFR-SEC-6}, \texttt{NFR-PRIV-7}).
\end{enumerate}
\par\medskip

\hypertarget{sec:b5-alternative-flows}{}
\addcontentsline{toc}{subsubsection}{B5. Alternative Flows}
\subsubsection*{B5. Alternative Flows}
\par\medskip

\hypertarget{sec:a1-invite-a-member-after-creation}{}
\paragraph*{A1. Invite a member after creation}
\par\medskip

At basic flow step 11, a member selects \textbf{Invite member}, chooses users in \texttt{InviteMemberPicker}, and sends invitations. The system shows “Invitation sent.” (\texttt{FR-33.4}, \texttt{AC-33.3}, \texttt{group\_detail.md}, recommendation).
\par\medskip

\hypertarget{sec:a2-remove-a-member}{}
\paragraph*{A2. Remove a member}
\par\medskip

At basic flow step 11, the authorized role selects \textbf{Remove from group}. \texttt{RemoveMemberConfirmDialog} asks “Remove \texttt{$<$name$>$} from \texttt{$<$group$>$}?” with Cancel/Remove (\texttt{FR-33.5}, \texttt{AC-33.4}).
\par\medskip

\hypertarget{sec:a3-leave-a-group}{}
\paragraph*{A3. Leave a group}
\par\medskip

The student selects \textbf{Leave group}. \texttt{LeaveGroupConfirmDialog} asks “Leave \texttt{$<$group$>$}?” with Cancel/Leave, and group content is no longer shown after confirmation (\texttt{FR-33.6}, \texttt{AC-33.5}, \texttt{feature\_flow.md}).
\par\medskip

\hypertarget{sec:a4-join-a-permitted-group}{}
\paragraph*{A4. Join a permitted group}
\par\medskip

From \texttt{My Group}, a permitted user selects \textbf{Join group}. \texttt{JoinGroupNotice} says “Joining depends on the group's privacy setting.” No approval or expiry rule is assumed (\texttt{feature\_flow.md}, \texttt{group\_detail.md}).
\par\medskip

\hypertarget{sec:a5-delete-a-group}{}
\paragraph*{A5. Delete a group}
\par\medskip

The \texttt{Owner} selects \textbf{Delete group}. \texttt{DeleteGroupConfirmDialog} says “Delete \texttt{$<$group$>$}? Messages and documents may be removed.” with Cancel/Delete (\texttt{overview.md}, \texttt{group\_detail.md}, recommendation).
\par\medskip

\hypertarget{sec:a6-report-inappropriate-group-behavior}{}
\paragraph*{A6. Report inappropriate group behavior}
\par\medskip

A member opens \texttt{ReportBehaviorDialog} from chat or the member list, provides category and context, and submits. The system shows “Report submitted for review.” without promising a moderator outcome (\texttt{FR-40.1}--\texttt{FR-40.4}, \texttt{AC-40.1}--\texttt{AC-40.4}).
\par\medskip

\hypertarget{sec:b6-exception-flows}{}
\addcontentsline{toc}{subsubsection}{B6. Exception Flows}
\subsubsection*{B6. Exception Flows}
\par\medskip

No group API error codes are specified; the following are source-backed UI states rather than backend error contracts.
\par\medskip

\hypertarget{sec:e1-access-denied}{}
\paragraph*{E1. Access denied}
\par\medskip

An unauthorized user opens a private group and sees “This group is private or you are not authorized to view it.” The system does not reveal group content (\texttt{NFR-SEC-4}, \texttt{NFR-SEC-6}, \texttt{group\_detail.md}).
\par\medskip

\hypertarget{sec:e2-invalid-uploaded-file}{}
\paragraph*{E2. Invalid uploaded file}
\par\medskip

The system shows “Unsupported file type” or “File does not meet size policy.” and does not accept the file (\texttt{FR-27.3}, \texttt{FR-15.3}, \texttt{NFR-SEC-5}). No numeric limit is asserted.
\par\medskip

\hypertarget{sec:e3-removal-right-conflict}{}
\paragraph*{E3. Removal-right conflict}
\par\medskip

The system uses \texttt{GroupMemberRoleBadge} to expose \texttt{Owner}, \texttt{Superuser}, and \texttt{Member}, does not silently merge the owner and Superuser rules, and flags the conflict for a product decision (\texttt{FR-33.5}, \texttt{AC-33.4}, \texttt{feature\_flow.md:145}).
\par\medskip

\hypertarget{sec:e4-invite-approval-or-expiry-expected}{}
\paragraph*{E4. Invite approval or expiry expected}
\par\medskip

No approval step, expiration period, or automatic acceptance is defined. The system confirms only “Invitation sent.” and does not invent a state (\texttt{group\_detail.md}).
\par\medskip

\hypertarget{sec:e5-group-fails-to-load}{}
\paragraph*{E5. Group fails to load}
\par\medskip

The system shows “Unable to load this group.” with Retry (\texttt{group\_detail.md}, recommendation); no group status code or error code is asserted.
\par\medskip

\hypertarget{sec:b7-business-rules}{}
\addcontentsline{toc}{subsubsection}{B7. Business Rules}
\subsubsection*{B7. Business Rules}
\par\medskip

\begin{enumerate}
\item A user can create a private or public friend group (\texttt{FR-33.1}, \texttt{AC-33.1}).
\item Group setup includes a name, privacy setting, and initial members (\texttt{FR-33.2}, \texttt{AC-33.1}, \texttt{feature\_flow.md}).
\item Members can send messages and share documents within the group (\texttt{FR-33.3}, \texttt{AC-33.2}).
\item Members can view the member list and invite new members (\texttt{FR-33.4}, \texttt{AC-33.3}).
\item The written requirement assigns member removal to the group owner (\texttt{FR-33.5}, \texttt{AC-33.4}); the flow separately names a Superuser kick, and the conflict is not resolved here.
\item Members can leave the group (\texttt{FR-33.6}, \texttt{AC-33.5}).
\item Uploaded files must pass type and size validation before acceptance (\texttt{FR-27.3}, \texttt{FR-15.3}, \texttt{NFR-SEC-5}).
\item Shared-resource access follows the sharer's privacy setting (\texttt{FR-27.5}, \texttt{FR-15.5}, \texttt{NFR-PRIV-5}, \texttt{NFR-PRIV-7}).
\item Private-group content is restricted to authorized users, with server-side authorization checks (\texttt{NFR-SEC-4}, \texttt{NFR-SEC-6}).
\item Inappropriate behavior can be reported and routed to moderation (\texttt{FR-40.1}--\texttt{FR-40.4}, \texttt{AC-40.1}--\texttt{AC-40.4}).
\end{enumerate}
\par\medskip

\hypertarget{sec:b8-special-requirements}{}
\addcontentsline{toc}{subsubsection}{B8. Special Requirements}
\subsubsection*{B8. Special Requirements}
\par\medskip

\begingroup
\small
\renewcommand{\baselinestretch}{1.0}\selectfont
\renewcommand{\arraystretch}{1.15}
\setlength{\tabcolsep}{6pt}
\begin{longtable}{|>{\raggedright\arraybackslash}p{0.10\textwidth}|>{\raggedright\arraybackslash}p{0.83\textwidth}|}
\hline
\textbf{ID} & \textbf{Requirement} \\ \hline
\endhead
NFR-SEC-1 & Authentication is required before accessing protected user resources. \\ \hline
NFR-SEC-3 & Sensitive communication must use TLS 1.2 or higher. \\ \hline
NFR-SEC-4 & Private-group content must be restricted to authorized users. \\ \hline
NFR-SEC-5 & Uploaded files must be validated for file type and file size before acceptance. \\ \hline
NFR-SEC-6 & Server-side authorization checks must prevent unauthorized access. \\ \hline
NFR-SEC-8 & The system must protect against common web vulnerabilities, including SQL injection, XSS, and CSRF. \\ \hline
NFR-PRIV-5 & Users must be able to control access to documents and learning resources they share. \\ \hline
NFR-PRIV-6 & Users must be able to control whether other users can send them messages or room invitations. \\ \hline
NFR-PRIV-7 & Privacy settings must apply consistently to member lists and search results. \\ \hline
NFR-USE-2 & The interface must support responsive layouts from 320 px to 1920 px. \\ \hline
NFR-USE-4 & Primary actions must use clear labels and recognizable controls. \\ \hline
NFR-USE-5 & The interface must be usable at screen widths of at least 320 px. \\ \hline
\end{longtable}
\endgroup
\par\medskip

\hypertarget{sec:b9-extension-points}{}
\addcontentsline{toc}{subsubsection}{B9. Extension Points}
\subsubsection*{B9. Extension Points}
\par\medskip

\begingroup
\small
\renewcommand{\baselinestretch}{1.0}\selectfont
\renewcommand{\arraystretch}{1.15}
\setlength{\tabcolsep}{6pt}
\begin{longtable}{|>{\raggedright\arraybackslash}p{0.35\textwidth}|>{\raggedright\arraybackslash}p{0.58\textwidth}|}
\hline
\textbf{Extension Point} & \textbf{Location in Flow of Events} \\ \hline
\endhead
Group creation & Basic flow steps 1-7 \\ \hline
Group collaboration & Basic flow steps 8-11 \\ \hline
Invite member & Alternative A1 \\ \hline
Remove member & Alternative A2; exception E3 \\ \hline
Leave group & Alternative A3 \\ \hline
Join group & Alternative A4 \\ \hline
Delete group & Alternative A5 \\ \hline
Report behavior & Alternative A6 \\ \hline
Access and file validation & Basic flow steps 10 and 12; exceptions E1-E2 \\ \hline
\end{longtable}
\endgroup
\par\medskip

\hypertarget{sec:b10-interface-and-data-notes}{}
\addcontentsline{toc}{subsubsection}{B10. Interface and Data Notes}
\subsubsection*{B10. Interface and Data Notes}
\par\medskip

\begin{itemize}
\item \textbf{Screen/flow path (confirmed from proposals):} \texttt{screen\_flow.md} gives \texttt{Home -$>$ Friend Group -$>$ Create Group}; \texttt{feature\_flow.md} gives \texttt{Friend -$>$ My Group -$>$ Group / Leave / Join}, \texttt{Group -$>$ Message / Document / List member / Leave}, and \texttt{Member -$>$ Invite / Kick by superuser}; \texttt{overview.md} gives \texttt{Create / Delete Group (Private | Public)} and \texttt{Friend / Group -$>$ Message / Material}.
\item \textbf{Recommended routes and components:} \texttt{/friends/groups/:groupId}, \texttt{/friends/groups/new}, \texttt{CreateGroupDialog}, \texttt{GroupNameInput}, \texttt{GroupPrivacySelect}, \texttt{BeginningMemberPicker}, \texttt{GroupDetailPage}, \texttt{GroupHeader}, \texttt{GroupMessagePanel}, \texttt{SharedDocumentPanel}, \texttt{GroupMemberPanel}, \texttt{InviteMemberButton}, \texttt{RemoveMemberButton}, and \texttt{LeaveGroupButton} are UI recommendations from \texttt{friend.md} and \texttt{group\_detail.md}.
\item \textbf{Exact role and privacy values:} roles are \texttt{Owner}, \texttt{Superuser}, and \texttt{Member}; privacy is \texttt{Private} or \texttt{Public}. No source defines how \texttt{Superuser} is granted, and no source defines a group capacity limit.
\item \textbf{Exact group fields:} \texttt{groupId}, \texttt{name}, \texttt{description}, \texttt{privacy}, \texttt{owner}, \texttt{memberCount}, \texttt{members}, \texttt{role}, \texttt{createdAt}, \texttt{lastActivityAt}, \texttt{messageId}, \texttt{messageText}, \texttt{messageSender}, \texttt{messageTime}, \texttt{documentId}, \texttt{documentName}, \texttt{documentType}, \texttt{documentSize}, \texttt{uploader}, \texttt{uploadedAt}, \texttt{inviteState}, and \texttt{reportContext} (\texttt{group\_detail.md}).
\item \textbf{Exact quoted UI text:} \textbf{Create group}, \textbf{Message the group}, \textbf{Send}, \textbf{Attach}, \textbf{Share material}, \textbf{Upload document}, \textbf{Invite member}, \textbf{Remove from group}, \textbf{Leave group}, \textbf{Delete group}, \textbf{Join group}, “Joining depends on the group's privacy setting.”, “Invitation sent.”, “Remove \texttt{$<$name$>$} from \texttt{$<$group$>$}?”, “Leave \texttt{$<$group$>$}?”, “Delete \texttt{$<$group$>$}? Messages and documents may be removed.”, “No messages yet.”, and “Report submitted for review.”
\item \textbf{Privacy notices:} “This group is private or you are not authorized to view it.”, “Private group content is visible to authorized members only.”, “Access follows the privacy set by the person who shared this file.”, and “This material may also appear in the Learning Resource Library.” (\texttt{group\_detail.md}).
\item \textbf{File validation messages:} “Unsupported file type” and “File does not meet size policy.”; the source explicitly states no numeric limit (\texttt{FR-27.3}, \texttt{FR-15.3}, \texttt{NFR-SEC-5}).
\item \textbf{Friend/profile boundary:} friend-request statuses \texttt{None}, \texttt{Pending}, \texttt{Accepted}, \texttt{Declined}, and \texttt{Blocked} and the 200-character request-note limit belong to UC-27/UC-28; group invitations have no specified status, approval, expiry, or note limit.
\item \textbf{API endpoints, statuses, and error codes:} Not specified in proposals. No group route, HTTP status, or error code is asserted, and generic banners such as “Unable to load this group.” are recommendations.
\item \textbf{Notification/privacy rules:} Invitation and social-activity notifications follow \texttt{FR-38.2}/\texttt{AC-38.3}; mute and block follow \texttt{FR-39.3}/\texttt{AC-39.1}; privacy enforcement follows \texttt{NFR-PRIV-5}, \texttt{NFR-PRIV-6}, and \texttt{NFR-PRIV-7}.
\end{itemize}
\par\medskip

\hypertarget{sec:b11-source-alignment-notes}{}
\addcontentsline{toc}{subsubsection}{B11. Source Alignment Notes}
\subsubsection*{B11. Source Alignment Notes}
\par\medskip

\begin{itemize}
\item There is no dedicated PA3 friend/group UC; UC-29 is derived from US-33, the friend/group flows, and UI recommendations in \texttt{friend.md}, \texttt{friend\_profile.md}, and \texttt{group\_detail.md}.
\item No friend/group API endpoints or error codes are defined in \texttt{API.md} or \texttt{response\_code.md}; group creation, messaging, documents, invitations, removal, leave, and delete have no specified backend contract.
\item The source conflicts on whether the group owner or Superuser removes members: \texttt{FR-33.5}/\texttt{AC-33.4} name the group owner, while \texttt{feature\_flow.md:145} says Superuser kicks. \texttt{GroupMemberRoleBadge} exposes \texttt{Owner}, \texttt{Superuser}, and \texttt{Member}, but no source defines how \texttt{Superuser} is assigned.
\item \texttt{AC-33.4} says “Owner Controls”, but \texttt{friend.md} and \texttt{group\_detail.md} recommend showing removal only to the owner/superuser; this document keeps the written requirement and flags the conflict rather than merging roles.
\item Group deletion is overview-derived from \texttt{overview.md}, not defined in \texttt{FR-33.1}--\texttt{FR-33.6}; owner scope and warning text are recommendations only.
\item \texttt{FR-33.3}/\texttt{AC-33.2} cover sending messages only; no group message edit, delete, reaction, or read-receipt rule is defined.
\item \texttt{FR-33.4}/\texttt{AC-33.3} cover invitations only; no invite approval step, expiration period, or automatic acceptance is asserted.
\item \texttt{group\_detail.md} explicitly excludes meeting-room rules; room creation, configuration, join requests, tools, and hand-wave behavior are not transferred to friend groups.
\item \texttt{settings.md} contributes communication and privacy controls for mute, block, and direct-message restriction; it does not add group API behavior.
\item Routes, component names, badges, dialogs, and quoted messages in \texttt{friend.md}, \texttt{friend\_profile.md}, and \texttt{group\_detail.md} are UI/UX recommendations derived from proposals, except where an exact FR or AC is cited above.
\end{itemize}
\par\medskip

\par\medskip
\hypertarget{sec:ui-prototype-specimens}{}
\addcontentsline{toc}{subsection}{UI Prototype Specimens}
\subsection*{UI Prototype Specimens}
\par\medskip
\begin{figure}[H]
\centering
\includegraphics[width=0.75\textwidth,height=0.42\textheight,keepaspectratio]{"../../prototype/friend_01_chat_aisha_patel.png"}
\par\vspace{2pt}
{\footnotesize \textbf{UI Prototype:} Friend: 01 Chat - Aisha Patel}
\end{figure}
\par\smallskip
\begin{figure}[H]
\centering
\includegraphics[width=0.75\textwidth,height=0.42\textheight,keepaspectratio]{"../../prototype/friend_01_chat_mia_chen.png"}
\par\vspace{2pt}
{\footnotesize \textbf{UI Prototype:} Friend: 01 Chat - Mia Chen}
\end{figure}
\par\smallskip
\begin{figure}[H]
\centering
\includegraphics[width=0.75\textwidth,height=0.42\textheight,keepaspectratio]{"../../prototype/friend_01_chat_noah_williams.png"}
\par\vspace{2pt}
{\footnotesize \textbf{UI Prototype:} Friend: 01 Chat - Noah Williams}
\end{figure}
\par\smallskip
\begin{figure}[H]
\centering
\includegraphics[width=0.75\textwidth,height=0.42\textheight,keepaspectratio]{"../../prototype/friend_01_chat_oliver_reed.png"}
\par\vspace{2pt}
{\footnotesize \textbf{UI Prototype:} Friend: 01 Chat - Oliver Reed}
\end{figure}
\par\smallskip
\begin{figure}[H]
\centering
\includegraphics[width=0.75\textwidth,height=0.42\textheight,keepaspectratio]{"../../prototype/friend_01b_no_conversation_selected.png"}
\par\vspace{2pt}
{\footnotesize \textbf{UI Prototype:} Friend: 01b No conversation selected}
\end{figure}
\par\smallskip
\begin{figure}[H]
\centering
\includegraphics[width=0.75\textwidth,height=0.42\textheight,keepaspectratio]{"../../prototype/friend_04_my_group.png"}
\par\vspace{2pt}
{\footnotesize \textbf{UI Prototype:} Friend: 04 My group}
\end{figure}
\par\smallskip
\begin{figure}[H]
\centering
\includegraphics[width=0.75\textwidth,height=0.42\textheight,keepaspectratio]{"../../prototype/friend_04_physics_lab_document.png"}
\par\vspace{2pt}
{\footnotesize \textbf{UI Prototype:} Friend: 04 Physics Lab - Document}
\end{figure}
\par\smallskip
\begin{figure}[H]
\centering
\includegraphics[width=0.75\textwidth,height=0.42\textheight,keepaspectratio]{"../../prototype/friend_04_physics_lab_list_member.png"}
\par\vspace{2pt}
{\footnotesize \textbf{UI Prototype:} Friend: 04 Physics Lab - List member}
\end{figure}
\par\smallskip
\begin{figure}[H]
\centering
\includegraphics[width=0.75\textwidth,height=0.42\textheight,keepaspectratio]{"../../prototype/friend_04_physics_lab_message.png"}
\par\vspace{2pt}
{\footnotesize \textbf{UI Prototype:} Friend: 04 Physics Lab - Message}
\end{figure}
\par\smallskip
\begin{figure}[H]
\centering
\includegraphics[width=0.75\textwidth,height=0.42\textheight,keepaspectratio]{"../../prototype/friend_04_study_circle_document.png"}
\par\vspace{2pt}
{\footnotesize \textbf{UI Prototype:} Friend: 04 Study Circle - Document}
\end{figure}
\par\smallskip
\begin{figure}[H]
\centering
\includegraphics[width=0.75\textwidth,height=0.42\textheight,keepaspectratio]{"../../prototype/friend_04_study_circle_list_member.png"}
\par\vspace{2pt}
{\footnotesize \textbf{UI Prototype:} Friend: 04 Study Circle - List member}
\end{figure}
\par\smallskip
\begin{figure}[H]
\centering
\includegraphics[width=0.75\textwidth,height=0.42\textheight,keepaspectratio]{"../../prototype/friend_04_study_circle_message.png"}
\par\vspace{2pt}
{\footnotesize \textbf{UI Prototype:} Friend: 04 Study Circle - Message}
\end{figure}
\par\smallskip
\begin{figure}[H]
\centering
\includegraphics[width=0.75\textwidth,height=0.42\textheight,keepaspectratio]{"../../prototype/friend_04a_calculus_crew_message.png"}
\par\vspace{2pt}
{\footnotesize \textbf{UI Prototype:} Friend: 04a Calculus Crew - Message}
\end{figure}
\par\smallskip
\begin{figure}[H]
\centering
\includegraphics[width=0.75\textwidth,height=0.42\textheight,keepaspectratio]{"../../prototype/friend_04b_calculus_crew_document.png"}
\par\vspace{2pt}
{\footnotesize \textbf{UI Prototype:} Friend: 04b Calculus Crew - Document}
\end{figure}
\par\smallskip
\begin{figure}[H]
\centering
\includegraphics[width=0.75\textwidth,height=0.42\textheight,keepaspectratio]{"../../prototype/friend_04c_calculus_crew_list_member.png"}
\par\vspace{2pt}
{\footnotesize \textbf{UI Prototype:} Friend: 04c Calculus Crew - List member}
\end{figure}
\par\smallskip
\begin{figure}[H]
\centering
\includegraphics[width=0.75\textwidth,height=0.42\textheight,keepaspectratio]{"../../prototype/friend_05_states.png"}
\par\vspace{2pt}
{\footnotesize \textbf{UI Prototype:} Friend: 05 States}
\end{figure}
\par\smallskip
\begin{figure}[H]
\centering
\includegraphics[width=0.75\textwidth,height=0.42\textheight,keepaspectratio]{"../../prototype/message_document_members.png"}
\par\vspace{2pt}
{\footnotesize \textbf{UI Prototype:} Message $\rightarrow$ Document $\rightarrow$ Members}
\end{figure}
\par\smallskip

\end{document}
